\documentclass[11pt]{article}
\usepackage[T1]{fontenc}
\usepackage[margin=1.04in,top=.91in,bottom=.92in]{geometry}
\usepackage{amsmath,amssymb,amsthm}
\usepackage[unicode,hidelinks]{hyperref}
\newtheorem{theorem}{Theorem}[section]
\newtheorem{lemma}[theorem]{Lemma}
\newtheorem{proposition}[theorem]{Proposition}
\newtheorem{corollary}[theorem]{Corollary}
\numberwithin{equation}{section}
\allowdisplaybreaks[2]
\setlength{\emergencystretch}{1.5em}
\newcommand{\PP}{\mathbb P}
\newcommand{\EE}{\mathbb E}
\newcommand{\Q}{\mathbb Q}
\newcommand{\one}{\mathbf 1}
\DeclareMathOperator{\rank}{rank}
\DeclareMathOperator{\spanop}{span}
\hypersetup{pdftitle={The normal order of the Erd\H{o}s--Hooley \unichar{"0394}-function},
            pdfauthor={Tseng}}
\title{The normal order of the Erd\H{o}s--Hooley $\Delta$-function}
\author{Tseng}
\date{}
\begin{document}
\maketitle

\begin{abstract}
Let $\Delta(n)$ be the largest number of divisors of $n$ whose
logarithms lie in an interval of length one. We prove that
$\Delta(n)=(\log\log n)^{\alpha+o(1)}$ on a set of natural density
one, with $0.35553850187<\alpha<0.3556265775100$. We identify
$\alpha$ with the equal-subset-sum exponent $\zeta_+$ of Ford, Green
and Koukoulopoulos. This proves the integer case of their Conjecture~2
and disproves the prediction $\zeta_+=\eta$ in their Conjecture~1,
where $\eta=0.35332277\ldots$ is the exponent of their binary
construction. The main step converts entropy constraints on
large-prime choices admitting small-prime completions into a
construction that can be repeated on disjoint ranges. Two variational formulas then give certified
bounds for $\alpha$ to arbitrary precision. The stated interval comes
from a construction with sixteen types of splits and a recursive
bound valid for every admissible tree.
\end{abstract}

\section{Introduction}\label{s:intro}

For a positive integer $n$, set
\[
 \Delta(n)=\max_{u\in\mathbb R}\#\{d\mid n:e^u<d\le e^{u+1}\}.
\]
We determine the growth exponent of this maximum for almost every
integer.

\begin{theorem}\label{t:main}
There is a real number $\alpha$ such that for every $\varepsilon>0$,
\[
 \lim_{x\to\infty}\frac1x\#\left\{16\le n\le x:
 \left|\frac{\log\Delta(n)}{\log\log\log n}-\alpha\right|
 >\varepsilon\right\}=0 .
\]
Equivalently, $\Delta(n)=(\log\log n)^{\alpha+o(1)}$ on a set of
natural density one. The exponent satisfies
\begin{equation}
             0.35553850187<\alpha<0.3556265775100 .
             \label{e:decimal}
\end{equation}
\end{theorem}

In 1948 Erd\H{o}s asked whether almost every integer has two
divisors $d_1<d_2<2d_1$ \cite{erdos}. In 1975
Erd\H{o}s and Nicolas studied the concentration of divisors in
dyadic intervals \cite{en75}. Their 1976 paper defines $F(n)$ as the
largest number of divisors of $n$ in an interval $(t/2,t]$ and studies
the integers at which $F$ attains a new record \cite{en76}. The two choices of
window give $F(n)\le\Delta(n)\le2F(n)$, since two dyadic intervals
cover a window of ratio $e$. In 1979 Hooley bounded the mean of
$\Delta$. He used this estimate to bound how closely $n^2\theta$
can approach a prescribed point modulo one for irrational $\theta$,
and to give a lower bound for the count of integers of the form
$a^2+b^4+c^4$ \cite[Theorems~4--5]{hooley}. Hall and Tenenbaum
studied both the mean and the normal order \cite{ht}.
Maier and Tenenbaum settled Erd\H{o}s's original density question
\cite{mt84}.

Two successive predictions for the normal-order exponent arose
from ways of building clusters of divisors. Maier and Tenenbaum
developed lower bounds by constructing close pairs on disjoint prime
supports \cite{mt85,mt09}. Taking one divisor from each of $m$ pairs
gives $2^m$ products, while the logarithmic widths of the pairs add.
In 2009 they obtained the lower exponent $0.33827\ldots$ and
conjectured that it was optimal \cite[Theorems~1.3--1.4 and the
conjecture following Theorem~1.4]{mt09}.
Ford, Green and Koukoulopoulos allowed a whole family of choices
on each range of primes. Their binary construction raised the lower
exponent to $\eta=0.35332277\ldots$ and disproved the Maier--Tenenbaum
prediction \cite[Corollary~1 and Theorem~3]{fgk}. They distinguished
two conjectures. Conjecture~1 predicts the value $\eta$ for the
exponent $\zeta_+$ of their equal-subset-sum model. Conjecture~2
predicts that $\zeta_+$ is sharp for integers, permutations and
polynomials over finite fields \cite{fgk}. For the integers we prove
$\alpha=\zeta_+$. The strict lower bound in \eqref{e:decimal} gives
$\zeta_+>\eta$, so the FGK binary construction also underestimates
the normal-order exponent. On the upper side, the exponent $\log2$ of Maier and
Tenenbaum \cite{mt09} was reduced to $0.61024\ldots$ by
de la Bret\`eche and Tenenbaum \cite{bt}.

A related question fixes the number $k$ of divisors and asks how
fast the interval containing them can shrink. Hou and Zhao
\cite[Theorem~2]{hz} and Mao and Song \cite[Theorem~1.2]{ms} determined
the supremum of the $a$ for which almost every $n$ has $k$ distinct
divisors with largest-to-smallest ratio at most $1+(\log n)^{-a}$.
It is $\beta_k/(1-\beta_k)$, where $\beta_k$ is the fixed-$k$
equal-sum threshold defined in Section~\ref{s:flags}. This settles
\cite[Conjecture~3]{fgk}. Mao and Song also proved that above the
threshold the integers possessing such a cluster have density zero.
Both papers identified the weak and strict entropy thresholds for
$\beta_k$. We use their refinement and perturbation method in
Section~\ref{s:flags}.

\subsection*{Idea of the proof}

A lower bound can be built one range of primes at a time. Once a
range supplies several choices with nearly equal logarithmic sums,
their products with the choices on a disjoint range multiply the
size of the cluster. Reversing this argument presents a difficulty.
Two close divisors may have widely separated large-prime parts.
Smaller primes can compensate for the difference, in many ways
that depend on the large-prime choices. An upper bound has to
allow all those completions.

Our comparison starts from the entropy constraints on a family
that admits completions. As successively smaller primes are read, their
membership patterns split the divisors into classes. The resulting
tree records when two divisors first differ, but distinct branches
can still satisfy a common linear constraint. We prove that the
entropy of the information shared by the branches suffices to remove
these constraints while preserving all split times. In the enlarged
spaces one can keep a fixed fraction of the rows and terminate the
system at a positive height. It now describes a construction on one
range, with the smaller ranges free for further repetitions. Both
the number of choices and the cutoff height change by at most
constant factors. Hence the upper count and the lower construction
give the same exponent.

Theorem~\ref{t:variation} makes these two descriptions precise.
Sections~\ref{s:open} and \ref{s:tree} prove the comparison, and
Section~\ref{s:arith} transfers it to the integers. Its quantitative
form gives the stopping rule in Section~\ref{s:eval}. The lower and
upper bounds in Sections~\ref{nl:section} and \ref{nu:section}
come from counting the leaves of one tree and bounding the growth
of every feasible tree, respectively.

\section{Two finite variational problems}\label{s:flags}

All logarithms are natural. For a finite random variable $X$ with
law $(p_x)$ write $H(X)=-\sum_x p_x\log p_x$, with $0\log0=0$.
We use $I(X;Y)=H(X)+H(Y)-H(X,Y)$ and relative entropy
$D(P\Vert Q)=\sum_x p_x\log(p_x/q_x)$, with value $+\infty$ if
$P$ charges a $Q$-null set. Conditional versions have their usual
meaning \cite{ct}. For a subspace $W$, $X\bmod W$ denotes the random
coset of $W$ containing $X$.

\subsection{Spaces, times and observations}

For a squarefree integer and an ordered list of $k$ divisors, give
each prime factor a column in $\{0,1\}^k$ recording divisibility in the
$k$ rows. Multiplying the column by $\log p$ and summing recovers
the vector of divisor logarithms. The constant vector changes
all row sums equally, so only their image modulo constants matters.
Ford, Green and Koukoulopoulos \cite[Sections~3--4]{fgk} record
the span of the columns as they are read in decreasing order.
The time $t=(\log\log p)/s$ has a useful normalization. In the
prime model a stage of length $c_j-c_{j+1}$ contains about
$s(c_j-c_{j+1})$ points on average, by
Lemma~\ref{l:prime-log-measure}.

Fix $k\ge2$. A flag is a chain of rational vector spaces
$V_0=\Q\one\subset V_1\subset\cdots\subset V_r\subseteq\Q^k$,
each obtained from the preceding space by adjoining one vector in
$\{0,1\}^k$. Its times satisfy $1\ge c_1\ge\cdots\ge c_r\ge b\ge0$,
with $c_{r+1}=b$. At times above $c_1$ the space is $V_0$.
In decreasing time, $V_j$ is born at $c_j$ and lasts until $c_{j+1}$.
A tie retains the order of the two births. In the $j$th stage choose
a law for a Boolean column $X_j$ in $V_j$. A test consists of
nested subspaces $W_0=V_0\subseteq W_1\subseteq\cdots\subseteq W_r$
with $W_j\subseteq V_j$. Put $d_j=\dim(V_j/W_j)$. The excess of
the test is
\begin{equation}
 \sum_{j=1}^r(c_j-c_{j+1})\{H(X_j\bmod W_j)-d_j\}-b d_r .
 \label{e:pay}
\end{equation}
A flag is weak if all excesses are nonnegative. It is strict if
every test other than $W_j=V_j$ for all $j$ has positive excess.
To interpret the two terms, prescribe the empirical law of the
columns on a stage. The multinomial count of their orders, as seen
modulo $W_j$, has logarithm governed by the number of columns times
$H(X_j\bmod W_j)$. On the other hand, each observed dimension
requires the corresponding row sums to match at the scale of its
birth. Summation by parts distributes this cost over the lifetime
of each dimension. A dimension still observed at $b$ incurs the
additional cost $b$. This explains the terminal term in \eqref{e:pay}.

It is convenient to allow intermediate stages on which the space
stays fixed. Their laws may be combined in a length-weighted mixture
by entropy concavity, testing first with a constant kernel throughout
the combined stage. Conversely, for a constant law, permitting extra
jumps of the kernel preserves weakness. Sample one time uniformly
in each positive-length stage and average \eqref{e:pay} for the
sampled kernels. They are nested and the terminal codimension can only decrease if the kernel grows at the
end of that stage. Zero-length intermediate
stages may be filled by intersecting the next kernel with the current
space.

The cells of $V$ are the classes of rows with identical evaluations
on $V$. A space separates the rows if every cell is a singleton.
Define $b_k$ as the supremum of the endpoints of strict separating
flags. Define $s_k$ as the greatest height at which all rows are
already distinguished in a weak flag with endpoint zero. Births at
the specified height count. The compactness proof for $s_k$ is given
below. Even for two rows these quantities differ. There is one
contrast, taking the three values $-1,0,1$. A birth at time one
and a uniform law on these three values give excess
$(1-b)\log3-1$ for the only proper test. This is the largest
possible excess, so $b_2=1-1/\log3$ and $s_2=1$.

\begin{theorem}\label{t:variation}
The exponent in Theorem~\ref{t:main} satisfies
\begin{equation}
 \alpha=\sup_{k\ge2}\frac{\log k}{\log(1/b_k)}
       =\inf_{\substack{k\ge2\\s_k<1}}
                  \frac{\log(k-1)}{\log(1/s_k)} .
 \label{e:alpha}
\end{equation}
For $k\ge10$ a single finite optimization gives
\begin{equation}
 \frac{\log\lfloor k/5\rfloor}{\log(8/s_k)}
       \le\alpha\le\frac{\log(k-1)}{\log(1/s_k)} .
 \label{e:singlek}
\end{equation}
Its width is $O(1/\log k)$ with explicit bound \eqref{e:width}.
For every rational $\varepsilon>0$, a finite algorithm returns
rationals $l\le\alpha\le u$ with $u-l<\varepsilon$.
\end{theorem}

\subsection{Finite tests and the arithmetic lower rate}

There are finitely many ambient chains because each new generator
is Boolean. To make the tests finite as well, it is essential to
record when two letters lie in the same observed coset. This is
the difference-vector reduction of Mao and Song
\cite[Lemma~2.2]{ms}, refining the subflag count in
\cite[Section~4]{fgk}. For stage alphabet
$\Omega_j=V_j\cap\{0,1\}^k$ and any test replace $W_j$ by
\[
 \spanop_{\Q}\!\left(\{\one\}\cup
       \bigcup_{i\le j}\{u-v:u,v\in\Omega_i,\ u-v\in W_i\}\right).
\]
This space is contained in $W_j$, is nested, and preserves every
coset partition of the relevant stage alphabet. Its rank cost
can only increase. A proper test remains proper. Enumerating all
subsets of the finite difference sets and taking their cumulative
spans therefore suffices for both weakness and strictness.

In particular all relevant inequalities are continuous in finitely
many clocks and probability vectors. Weak feasibility is closed.
There are finitely many choices of the first separating space, so
compactness of the ordered clocks and the probability simplexes
proves attainment of $s_k$.

The passage from weak to strict inequalities at positive endpoints
is due to Hou and Zhao \cite[Theorem~4]{hz} and Mao and Song
\cite[Theorem~3.4]{ms}. Their method refines a flag into rank-one
jumps and mixes its laws with a system having strict slack. We use
an auxiliary system with a positive endpoint even when the original
one ends at zero. On the same rank-one Boolean chain, with successive
pivots $u_j$, take births $g^{j-1}$, endpoint $g^r$ and stage
letters $0,u_j,\one-u_j$ with equal probabilities, where
$(1-g)\log3>1$ and $0<g<1$. Use the ordered pivot basis modulo
constants. The possible last nonzero indices in a nested test
space form a nondecreasing set, whose cardinality is its dimension
modulo constants by elimination. For an index $j$ missing when
born, its entire rank cost including the endpoint is at most
$g^{j-1}$, whereas on its birth stage the three cosets are
distinct, earning $(1-g)g^{j-1}\log3$. Indices present at birth
cost nothing. A proper test misses some index. Thus this auxiliary
clock is strict. Mix the two clocks with a small positive weight
on the auxiliary one, using mixtures of the laws in proportion to
their weighted lengths on each stage. Entropy concavity and the
linear rank costs give strictness. All positive stages can also
be given full letter support by a sufficiently small perturbation
because there are only finitely many test profiles. The argument
applies to a zero endpoint as well. A positive weak flag conversely
extends weakly to zero using the zero letter, since the tail cost
is at most its previous terminal debit.

In the logarithmic random set $\mathcal A$ of Ford, Green and
Koukoulopoulos, each positive integer $m$ occurs independently with
probability $1/m$. Their $\beta_k$ is the supremum of the $b$ for
which $\mathcal A\cap(D^b,D]$ contains $k$ distinct subsets with
equal sums with probability tending to one as $D\to\infty$.
To identify this threshold with $b_k$, write $\widetilde{\gamma}_k,\gamma_k$ for the strict
and weak endpoint suprema of Ford, Green and Koukoulopoulos.
Their entropy criterion gives
$\widetilde{\gamma}_k\le\beta_k\le\gamma_k$
\cite[Theorem~7]{fgk}. Mao and Song give the corrected subflag and
residual-sum counts and verify this comparison
\cite[Lemmas~2.2--2.3 and Proposition~2.4]{ms}. In the energy
difference for a subflag, summation by parts gives
\[
 \sum_{j=1}^r c_j(d_j-d_{j-1})
 =\sum_{j=1}^r(c_j-c_{j+1})d_j+bd_r,\qquad d_0=0,
\]
which is precisely the dimension cost in \eqref{e:pay}.
The ambient spaces in the finite entropy criterion are spanned by
Boolean columns and its terminal nondegeneracy condition is exactly
row separation. Dividing all times by $c_1$ normalizes the first birth to one,
scales every excess by $1/c_1$ and can only increase the endpoint.
A jump by several Boolean generators can be refined by adjoining
them successively at the same time, retaining the original law on
the final stage of the jump. Compression to its last kernel preserves
the excess. The strict perturbation therefore identifies $b_k$ with $\beta_k$,
the common positive threshold in \cite[Theorem~4]{hz} and \cite[Theorem~1.5]{ms}. The height
$s_k$ permits the flag to continue after all rows have separated.
This continuation is what makes it useful for counting completions.

Put $A=\sup_{k\ge2}\log k/\log(1/b_k)$. We claim that $A$ equals the lower rate of \cite{fgk}, namely
\[
       \zeta_+=\limsup_{k\to\infty}
                    \frac{\log k}{\log(1/\beta_k)} .
\]
Fix $0<b<\beta_k$.
On the $m$ disjoint intervals $(D^{b^{j+1}},D^{b^j}]$,
$0\le j<m$, the equal-sum event for $k$ choices holds simultaneously
with probability tending to one for each fixed $m$. Unions give
$k^m$ distinct equal-sum subsets. Thus $\beta_{k^m}\ge b^m$, and
the limsup over the number of choices equals the supremum defining
$A$. The divisor transfer of \cite[Theorem~3]{fgk} gives, for every
$\varepsilon>0$,
\begin{equation}
  \Delta(n)\ge(\log\log n)^{A-\varepsilon}
  \qquad\text{for almost every }n .
  \label{e:arithlower}
\end{equation}
Section~\ref{s:arith} and Appendix~\ref{s:realization} give a direct
proof of \eqref{e:arithlower} on the same prime process used for
the upper bound. In particular the proof of the normal order can
be read independently of the random-set transfer.

\section{From separated rows to a positive endpoint}\label{s:open}

A new Boolean column splits each
affected cell into two children. This produces a rooted binary
tree, with ties allowed between splits. Restriction to a subset of
rows preserves weakness. The inverse images of a test under the
stagewise restriction maps have the same observed cosets and ranks.
Throughout a stage use one coordinate per current cell. On a node $v$ write $V_v$ for the restriction of the given space.
After the split at $v$, $V_v$ contains the two separate child
indicators, since the splitting column restricted to $v$ and its
complement are those indicators. The full space on the current
cells consists of all rational functions on those cells and contains
every Boolean word. Its chain, with the same
separation times, is called the \emph{full partition flag}. Ties may
be resolved at zero duration by adjoining the split indicators
one at a time. After the last of these births the space is the
full space on the refined partition.

For a block $v$, its raw word is the Boolean word on its current
cells. Its unrooted word is its coset modulo the constants
on $v$. The image determines each nonconstant raw word because the only
distinct Boolean words differing by a nonzero constant are the
two constant words. We call their common image \emph{vacant}.
At a split the two unrooted child words have respectively one,
two or three possible extensions to an unrooted parent according
as zero, one or two children are vacant. In the last case the
extensions are the vacant parent and the two opposite assignments.

\begin{theorem}[Restoration of the full partition space]\label{t:restore}
Every finite weak flag with endpoint zero has a weak full partition
flag on its tree with the same separation times.
\end{theorem}

Consider four rows and the successive columns
\[
             u=(0,0,1,1),\qquad v=(0,1,0,1).
\]
The first divides the rows into two pairs. The second distinguishes
both members of each pair. All four rows are now distinct, although
the space spanned by $\one,u,v$ has dimension three. Its two sibling
contrasts obey
\begin{equation}
                  x_2-x_1=x_4-x_3 . \label{e:fourrows}
\end{equation}
The full partition space is $\Q^4$, in which these contrasts may vary
separately. Row separation is therefore visible from the tree, while
the rank available at the endpoint also depends on the linear
relations between branches. Theorem~\ref{t:restore} supplies the
full rank needed in the next argument. Its proof is given in
Section~\ref{s:tree}.

For $0<c\le1$, let $M(c)$ be the greatest number of cells distinguished by height
$c$ in weak zero-endpoint flags, allowing arbitrary initial arity.
It is finite for each $c>0$. To see this take any internal node
at height $S$ with an internal child at $s$. Select two rows
across the latter split and one on the opposite side of the former.
Test all their contrasts throughout time. They have ranks one on
$(s,S)$ and two on $(0,s)$, and their respective entropies are
at most $\log3$ and $\log7$, by counting unrooted words.
The early stage can earn a surplus since $\log3>1$, while two
contrasts among three rows force a deficit since $\log7<2$.
In particular weakness forces a gap between the two split heights.
\[
 0\le(S-s)(\log3-1)+s(\log7-2),\qquad
 s\le\gamma S,\quad
 \gamma=\frac{\log3-1}{\log3-1+2-\log7}<1 .
\]
At most $1+\log(1/c)/\log(1/\gamma)$ splits of height at least
$c$ occur on a path. In particular $M(c)\le2c^{-a_0}$ with
$a_0=\log2/\log(1/\gamma)$. By restriction to representatives
$M(c)\ge k$ is equivalent to feasibility on just $k$ rows,
distinguished at $c$. By attainment of $s_k$ this is equivalent to $c\le s_k$.

Stopping a flag at positive height leaves a terminal rank cost.
A test with negative excess identifies the rows responsible for
this loss. Its kernel tells us which coordinate evaluations can
be kept, and its codimension is bounded by the entropy available
in the discarded tail.

\begin{lemma}\label{l:opening}
If $M(c)=R\ge10$, then $b_{\lfloor R/5\rfloor}\ge c/8$.
\end{lemma}
\begin{proof}
Restore the full flag and keep one row per cell present at $c$.
The restricted flag is weak and already has space $\Q^R$ at $c$.
One can use a height $s<c$ arbitrarily close to $c$ to be inside
a stage. For prefix tests ending at $s$, put
$G(W)=\int_s^1[H(X\bmod W)-\dim(V/W)]\,dt$. Every $G(W)$
is nonnegative by closing the kernel to the full space for the
rest of the clock. Set $b=s/8$ and minimize
$E(W)=G(W)-b\dim(\Q^R/W(s))$ over prefix tests on the finite
list of stage boundaries augmented by $s$. The finite test argument
in Section~\ref{s:flags} gives a minimizer. Shifting the prefix heights down by $s-b$ and filling the top
by constants gives precisely this positive-endpoint payoff. If the minimum is nonnegative
we keep every row. Otherwise let a minimizer $W_*$ have terminal
codimension $e$. Use its prefix on the old flag and keep its
kernel fixed until time zero. The tail observation has entropy
at most $R\log2$ at each time, so
\[
       0\le G(W_*)+s(R\log2-e)<be+s(R\log2-e).
\]
Thus $e<(8/7)R\log2<4R/5$. Keep $R-e$ independent coordinate
evaluations on $W_*(s)$, so their projection maps it onto the
whole new terminal space. If $U$ is the pullback of any test on these
rows, then $U+W_*$ is full at $s$. The entropy of a linear
quotient is submodular in its kernel. Observing modulo the
intersection is equivalent to observing both cosets, and the coset
of the sum is a function of each. Dimensions satisfy the corresponding
identity with equality. Integration therefore gives
\[
  E(U)\ge E(U+W_*)+E(U\cap W_*)-E(W_*)\ge G(U+W_*)\ge0.
\]
The projected flag is weak at positive endpoint $b$. Restrict
further if necessary, merge repeated spaces, and use the strict
perturbation above. Let $s\uparrow c$.
\end{proof}

The lower rate $A$ is positive and finite. Positivity follows from the two-row calculation. Restriction gives $b_k\le b_2$,
using weak restriction and strict approximation if necessary.
Filling to zero gives $b_k\le(2/k)^{1/a_0}$. This bounds the
ratios defining $A$ uniformly for large $k$, while $b_2<1$ covers
the remaining terms. From Lemma~\ref{l:opening} and $b_m\le m^{-1/A}$,
\begin{equation}
                    M(c)\le 10+5(8/c)^A . \label{e:growth}
\end{equation}
For $k=M(c)+1$ we have $s_k\le c$. It follows from \eqref{e:growth},
letting $c$ decrease to zero, that
\[
 \inf_{k:s_k<1}\frac{\log(k-1)}{\log(1/s_k)}
       \ \le\ \liminf_{c\downarrow0}\frac{\log M(c)}{\log(1/c)}
       \ \le\ A .
\]
Write $U_*$ for the infimum on the left. By \eqref{e:arithlower}
the lower rate $A$ holds on the integers. Section~\ref{s:arith} will
give the upper rate $U_*$ there. Since the two bounds hold in
natural density, $A\le U_*$. The inequality just proved then yields
$A=U_*$ and identifies the exponent. We turn to the two ingredients of this comparison, preservation of
weakness under restoration and the arithmetic bound on continuations.

\section{The full partition theorem}\label{s:tree}

In \eqref{e:fourrows} take the Boolean word uniformly from
$0,v,\one-v$. Each sibling contrast is uniform on $\{-1,0,1\}$,
but the pair carries only $\log3$ units of entropy since its
coordinates agree. Under the product of the two contrast laws
the pair has entropy $2\log3$ and its linear rank rises from one
to two. The common observation is the single contrast seen on
both sides. In general its entropy can be small at a particular
stage. The original weak flag bounds its entropy integrated over
time, because the common observation itself defines a test.

There is a further constraint on this idea. A test may see only
the joint word on a parent. If both child words are constant,
their two bits can agree or disagree, and the information in their
relative orientation affects the parent entropy. The replacement
laws must accommodate this possibility for every test at once.
Lemma~\ref{l:cutmin} reduces the tests to cuts of the tree.
Lemma~\ref{l:gibbs} constructs laws for nested cuts with a
bound of $2/3$ on each contrast atom. The comparison
\eqref{e:comparison} then recovers the entropy of the common
observations while accounting for the possible constant words.
Integration converts this recovery into the required rank
inequalities.

We prove the theorem on a fixed finite list of stages. Include the
zero-length stages needed for the two chains. At positive stages
the two chains have identical cells. When an old pivot splits
several cells at once the intermediate full spaces have duration
zero. Pulled-back kernels on the original positive stages extend
over omitted zero stages by intersecting the next kernel with
the intermediate ambient space (use the full kernel after the
last positive stage if necessary). All these extensions are
nested. The proof produces one law per positive-length stage.

\subsection{Reduction to cuts of the tree}

An evolving cut starts with the root. A cut member may be replaced
by its two children at a time when it has exactly two distinguished
cells. Its kernel is the space of words constant separately on
the current cut members. Its observation is the tuple of their
unrooted words. On the four rows in \eqref{e:fourrows}, cutting
into the two pairs observes $(x_2-x_1,x_4-x_3)$ and has kernel
$\{(a,a,b,b):a,b\in\Q\}$. Use only listed boundaries for these
decisions. Ford, Green and Koukoulopoulos use submodularity and
symmetries of a fixed flag to reduce the class of entropy tests
\cite[Section~8]{fgk}. Here the symmetries are sign changes on
individual subtrees. They force a minimizing kernel to decompose
as soon as the corresponding branches can be separated.
At each positive stage let $\mathcal L$ consist of laws with the
following two properties. Each pair of distinct cells agrees with
probability at most $1/2$. Complementing any entire nonconstant
rooted subtree, leaving it alone if it is constant, preserves
the raw law. These constraints define a compact convex set.

\begin{lemma}\label{l:cutmin}
For stage laws in $\mathcal L$, the minimum of the payoff over all
nested rational tests equals its minimum over evolving cuts.
\end{lemma}
\begin{proof}
On a block of $N\ge3$ current cells let $K$ contain the constants
and have dimension $r$. Then
\begin{equation}
 H(X\bmod\one\mid X\bmod K)\le r-1 . \label{e:fiber}
\end{equation}
Here and below $\bmod\one$ denotes the quotient by the constant line.
For $r=1$ this is immediate. In an affine $K$-fiber its Boolean
points are determined by $r$ independent coordinates. If all
$2^r$ patterns are attained, the two extreme patterns differ by
$\one$ on those coordinates and hence globally. Thus there are at
most $2^r-1$ distinct unrooted possibilities, whose logarithm
for $r\ge3$ is at most $r-1$. For $r=2$, put $K=\spanop(\one,v)$.
Three distinct unrooted possibilities have three distinct slopes
along $v$. If $v$ takes at least three values, each pairwise
nonconstant difference, having values in $\{-1,0,1\}$, has
oscillation two. Three distinct real slopes cannot have identical
pairwise distance. If $v$ has just two levels, three possibilities
require a word constant within both levels, since a nonconstant
level fixes its raw assignment and the other has at most two choices.
One level contains a fixed pair of cells. The event that the fiber
has three possibilities therefore has probability at most $1/2$.
The conditional entropy is at most $(\log2+\log3)/2<1$.

By submodularity, as in the proof of Lemma~\ref{l:opening}, the sum
of the excesses of two tests dominates the sum of the excesses of
their intersection and sum. A minimizer exists because
the dimensions and the induced partitions on all finite alphabets
take finitely many profiles. Choose one maximizing the sum of
kernel dimensions on all listed stages. By the preceding inequality
it contains every other minimizer.

Starting with the root, keep this test decomposed across the
present cut. Close a cut block as soon as its test kernel is the
whole space on its two current cells. At this closing stage
the kernel is a direct sum across the present cut and contains
both child indicators. It is therefore invariant under negation
on either child. Leave all kernels through this stage untouched
and apply one such negation to every later kernel. Invariance at
the closing stage gives nesting across the boundary. The payoff
is preserved because complement acts as negation modulo the
child constants and the vacant word maps to zero. Maximality
makes the chosen later kernels invariant
under each separate negation. Subtracting a negated vector and
dividing by two shows these kernels split as direct sums across
the new cut. This inductively justifies all the closures, also when
several occur together. Relative to observing all current unrooted
cut words, the entropy lost by this test is at most the sum on its
blocks of the left sides of \eqref{e:fiber}, by the conditional
entropy chain rule. A remaining block of size one or two has only
the constant kernel. The rank saved thus covers the lost entropy.
The cut has at most the minimizing payoff and is itself admissible.
\end{proof}

One law must pass all the cuts. Optimizing the law separately
for each cut would interchange the quantifiers in that requirement.
A separating hyperplane instead allows us to fix a weighted
collection of cuts and optimize a single law for their combined
payoff. Submodularity then orders the cuts into a chain, meaning
simultaneous inclusion of their kernels at every time.
For laws restricted to $\mathcal L$ on every positive stage,
the maximum minimum cut payoff equals
\begin{equation}
 \min_{\substack{C_1,\ldots,C_h\ {\rm a\ chain}\\
                       a_i\ge0,\ \sum_i a_i=1}}
     \ \max_{Q\in\mathcal L}\sum_i a_i\,g_{C_i}(Q),
 \label{e:minimax}
\end{equation}
where $g$ is the whole payoff. Indeed the downward closure of
the attainable vectors of cut payoffs is convex by entropy
concavity and closed by compactness. Separation from a strictly
larger constant vector gives a nonnegative supporting normal of
sum one, proving finite minimax. Two cuts can be replaced by
their meet and join. A split closes at the later or earlier of the
two closure times, including infinite times. The resulting cuts
are legal by the ancestry condition. In the basis of the constant and one child indicator per split,
each kernel is spanned by the previously closed splits and the
constant. The new kernels are therefore the intersection and sum. Thus this replacement decreases
the sum of payoffs for every law. Replicate rational weights into
equal weights and sort using meet and join as comparators in a
fixed sorting network (for example, repeated adjacent sweeps).
Sorting the zero-one indicator of every closure at every stage
puts the resulting whole tests in a chain. Approximation of the real
weights gives \eqref{e:minimax}, since all payoffs are uniformly
bounded and the list of cut chains is finite.

\subsection{Entropy maximizers on a full subtree}

At one stage consider nested forest observations of a full
Boolean word, with nonnegative weights, on a given subtree.
The objective is the weighted sum of their entropies, plus
$\theta$ times the probability of whole-subtree vacancy. Assume
the full unrooted word occurs with weight at least $w>0$ and
$0\le\theta\le w\log2$. Give the two lifts of a vacant root
equal probabilities to obtain a raw law. The reward accounts for
the bit that distinguishes these two lifts. Recursive optimization
over successively finer observations also occurs in the coset
trees of \cite[Section~7.2]{fgk}. Here the observations arise
from cuts of the row tree, and the same recursion must control
both the vacancy reward and the concentration of linear contrasts.

\begin{lemma}\label{l:gibbs}
This objective has an optimizer $Q$ in $\mathcal L$. It has full
support on the unrooted alphabet. For any nonzero scalar contrast
$L$ on the cells, meaning a linear functional annihilating constants,
\begin{equation}
                 \max_z Q(L=z)\le 2/3. \label{e:atom}
\end{equation}
If $G$ denotes the objective and $\mu$ any competing law, then
\[
 G(Q)-G(\mu)=\sum_i a_i
           D(\mu_{O_i}\Vert Q_{O_i}),
\]
where $O_i$ are the forest observations. When both raw laws
are complement symmetric the full unrooted divergence equals
the raw divergence. With zero vacancy reward and no required
full-word weight, an optimizer in $\mathcal L$ still exists.
\end{lemma}
\begin{proof}
Starting from individual vacant cells, reveal joins in the order
in which the nested views require them, children before parents.
A view consists precisely of all joins in its forest blocks.
Refine simultaneous reveals arbitrarily in this order. Number the
joins $1,\ldots,m$ and put $w_k$ equal to the total weight of the
views that see join $k$. The chain rule expresses the objective
as the sum of the conditional entropy at each join times $w_k$,
and the final vacancy reward. Here $w_1\ge\cdots\ge w_m\ge w$.
Let $y_k$ be the entire state after $k$ joins. Initialize
$F_m=\theta\one_{\rm root\ vacant}$ and compute backwards
\[
 F_{k-1}(y)=w_k\log\sum_{x:\ x\ {\rm extends}\ y}
                                    \exp(F_k(x)/w_k).
\]
Conditionally on $y$, assign probabilities proportional to
$\exp(F_k(x)/w_k)$. The identity for entropy minus relative
entropy verifies optimality and gives positive probabilities.
If the tree has a single cell the conclusions hold directly.

Each $F_k$ depends only on vacancy indicators of its present
blocks. Its increment on changing one block from nonvacant to
vacant is between zero and $w_{k+1}\log2$ for $k<m$.
The terminal bound uses $\theta\le w_m\log2$. In backward
induction an unchanged block inherits the bound by monotonicity
and constant translation of log-sum-exp. For either child of the
current join the increment is $w_k\log2$ if its sibling is
nonvacant. With a vacant sibling it is
$w_k\log[(2+\exp(\delta/w_k))/2]$, where $\delta$ is the
continuation increment for the parent and
$0\le\delta\le w_k\log2$. The induction also applies if one includes formal
nonvacant singleton states when computing the increments. Thus at a double vacancy the conditional chance of a
vacant parent is between $1/3$ and $1/2$. All available sign
choices are fair and vacancy decisions depend only on previous
vacancies. Given the whole vacancy pattern the compatible raw
assignments are equiprobable because the number and sites of the
fair choices are fixed and the signs can be recovered from the
assignment. Complementing a nonconstant rooted block leaves
this pattern unchanged. Two cells in different children of
their least common ancestor agree with conditional probability
$1/2$ if one child is nonvacant, by its complement symmetry.
When both children are vacant their agreement probability at that
join is at most $1/2$. This proves the two properties of $\mathcal L$.

For the scalar estimate, the log probability of a raw word,
up to a constant, is
\[
   \sum_{k=1}^{m-1}(1/w_k-1/w_{k+1})F_k(y_k)
       +(\theta/w_m-\log2)\one_{\rm root\ vacant}.
\]
This follows by multiplying the transition probabilities, including
the fair sign on a vacant root. All coefficients are nonpositive.
Fix one leaf $i$. If its sibling block $S$ is nonvacant, flipping
that leaf leaves all vacancy states unchanged. If $S$ is constant,
flipping the leaf away from that same bit can only remove vacancies.
Its raw probability can therefore only increase. Pairing words by
the flip shows that the total variation distance between $Q$ and
its flipped law is
\[
 Q(S\ {\rm vacant})-2Q({\rm parent\ vacant})
                \le Q(S\ {\rm vacant})/3\le1/3 .
\]
The first inequality uses the lower bound $1/3$ at the join of
the leaf with $S$. Flip a leaf with nonzero coefficient in $L$.
An atom event and its flipped event are disjoint, giving
$2Q(L=z)\le 1+1/3$.

At an interior maximum on the simplex of full unrooted words the
first variation vanishes on each zero-sum direction. Expanding the
entropies about that point gives the divergence identity, since
the vacancy reward is linear. A symmetric raw law is the same
fair extension at a vacant word, proving the divergence assertion.
With zero full-word weight and zero reward, add $\epsilon>0$ to
the full weight and pass to an optimizing subsequential limit.
\end{proof}

\subsection{The common observation and a reconstruction of raw words}

The contrast in \eqref{e:fourrows} has an intrinsic description
which also applies when the children have more than two cells.
For two linear observations on the given space $V$, with kernels
$T_L,T_R$, their common observation is $J=X\bmod(T_L+T_R)$.
It is a function of either observation. Its rank $j$ satisfies
\[
  \rank(Y_L,Y_R)=\rank Y_L+\rank Y_R-j,\qquad
  I(Y_L;Y_R)=H(J)+I(Y_L;Y_R\mid J).
\]
At an active split, the raw child observations and the unrooted
child observations give the same common quotient. Indeed any pair
of child functionals agreeing on $V_v$ must annihilate each child's
constant line, by evaluating on its separate indicator. Therefore
the common observation is zero if either child is vacant.

We first symmetrize the old raw law by the global complement at
every stage. This preserves weakness by concavity because it
permutes the cosets in each test. All restrictions are then
complement symmetric. For a symmetric raw word $Z$ and its
unrooted image $\bar Z$,
\begin{equation}
             H(Z)=H(\bar Z)+(\log2)\PP(Z\text{ vacant}).
 \label{e:raw}
\end{equation}
Here vacancy of the raw word means that it is constant.

\begin{lemma}\label{l:join}
Let $Z_L,Z_R$ be the old raw child words, and let their laws
already have symmetric replacements with vacancy probabilities
no greater than before. The two replacements can be coupled so
the parent vacancy probability does not increase and their mutual
information is at most $I(Z_L;Z_R\mid J)$ for the old common
observation $J$.
\end{lemma}
\begin{proof}
Write $q=\PP(J=0)$. Draw two independent labels $J_L,J_R$ with
the old law of $J$. When both labels are zero, draw the old pair
conditional on $J=0$. For every other label pair, draw the two
words independently from their respective old conditional marginals.
Each word has its old marginal, and its conditional marginal depends
only on its own label. By the chain rule for relative entropy, the
mutual information of the two words, each accompanied by its label, is
\[
 q^2 I(Z_L;Z_R\mid J=0)\le I(Z_L;Z_R\mid J).
\]
The left side is zero if $q=0$. A constant child forces its label
to be zero. Thus the probability that both children are constant
with the same bit is $q$ times its old value. The pair is
complement symmetric, since complement sends both labels to their
negatives and the old joint law is complement symmetric.

On each side retain the three states constant zero, constant one
and nonconstant. Independently thin the vacancies to the prescribed
replacement probabilities, keeping the bit of every retained vacancy.
Given the resulting states, draw words independently from the
replacement marginals conditional on their respective states.
These separate channels preserve the prescribed marginals and joint
complement symmetry and cannot increase mutual information. Thinning
cannot create a pair of equal constant children.
\end{proof}

At the remaining join, two independent reference words must be
compared with words whose common labels agree. The event that the
whole block is constant lies in the zero label on both sides.
The following estimate isolates the contribution of precisely that
event to the entropy comparison.

\begin{lemma}\label{l:tilt}
On a finite product let $Q=\bigotimes_i Q_i$. Two distinct factors
have functions $J_1,J_2$ taking values in the same space. Write
$q_{iz}=Q_i(J_i=z)$ for these two factors,
$h_0=\sqrt{q_{10}q_{20}}$, and let $E$ be an event of product mass
$b$ contained in $\{J_1=J_2=0\}$. For every law $P$ supported
on $J_1=J_2=:J$ and every $d>0$,
\[
 D(P\Vert Q)-H_P(J)+dP(E)
           \ \ge\ (1-e^{-d})b/h_0
\]
when $h_0>0$. When $h_0=0$, the left side is nonnegative and $b=0$.
\end{lemma}
\begin{proof}
Given both labels $z$, compare the conditional law of all factors
with the product reference conditioned on those labels. The
decomposition of $D(P\Vert Q)-H_P(J)$ is the mean of these
conditional divergences plus
$2\sum_z p_z\log(p_z/\sqrt{q_{1z}q_{2z}})$.
On $z=0$ exponentially tilt the conditional reference by
$e^{-d\one_E}$, giving the normalizing constant
$1-(1-e^{-d})b/h_0^2$. The conditional divergence to the
tilted law is nonnegative. On all other labels use the
unmodified reference. Minimizing the resulting expression in
$(p_z)$ gives
\[
 -2\log\left(\sum_{z\ne0}\sqrt{q_{1z}q_{2z}}
                  +\sqrt{h_0^2-(1-e^{-d})b}\right).
\]
The sum before reducing the zero summand is at most one. Its
reduction is at least $(1-e^{-d})b/(2h_0)$. Use
$-\log(1-t)\ge t$. Omit null labels. The inequality holds directly for an infinite
divergence.
\end{proof}

\subsection{Paying for a common observation}

Fix a stage and cuts $C_i$ with weights $a_i$ from
\eqref{e:minimax}. On an old node $v$, define $Y_{v,i}$ as
the unrooted word on $v$ if a member of $C_i$ contains $v$,
and otherwise as the tuple of unrooted words on the members
inside $v$. Let $J_{v,i}$ be the common observation of the
two children's respective $Y$ on the old space $V_v$, with rank
$j_{v,i}$. Write $Y_i$ for the root observation. We claim that
a stage law in $\mathcal L$ on the full space satisfies
\begin{equation}
 \sum_i a_i H_Q(Y_i)\ \ge\
       \sum_i a_i\left(H_P(Y_i)+\sum_{v\ {\rm split}}H_P(J_{v,i})\right).
 \label{e:comparison}
\end{equation}
Here $P$ is the given symmetric law and $Q$ is the law to be constructed.

Take as pieces the maximal rooted subtrees $S$ whose old
restricted space is the full space on their current cells.
Their maximal roots partition the cells, since a singleton
has full space and fullness is inherited by restrictions.
A cut member $D$ that meets more than one piece is a union of
whole pieces. Call such a block mixed. Each piece lies in at
most one mixed block of a given cut. For each $S$ let $A_{S,i}$
be its induced unrooted forest observation. This is the full
unrooted word on $S$ when it lies in a mixed block.
Let $w_S$ be the sum of the weights of those mixed occurrences.
Independently on the pieces take the optimizers of
\[
  G_S(\mu)=\sum_i a_i H_\mu(A_{S,i})
                   +w_S(\log2)\mu(S\text{ vacant})
\]
in Lemma~\ref{l:gibbs}, including fair orientations on vacancy.
This law $Q$ on all raw words belongs to $\mathcal L$. For a
subtree within a piece use the lemma. For a subtree that is a
union of pieces, complement of all those pieces preserves their
product law, and its nonconstancy is an invariant event. For
cells in different pieces agreement has probability $1/2$ by
independence and symmetry.

Fix one cut temporarily. Form terminals of a smaller tree from
all mixed cut blocks $D$ and all whole pieces $S$ outside them.
On each such $S$ the datum is the joint list $A_{S,i}$, even if
several cut members lie inside it. On a mixed terminal it is
$Y_D$, the whole unrooted word. Above these terminals the
observation at every node is just the pair of its child
observations. Its joint entropy plus the associated $H(J_{v,i})$
is at most the sum of their entropies. All common ranks inside
a full piece vanish, including for partial forest observations.
Consequently the contribution on the right of \eqref{e:comparison}
at this cut is bounded by the sum of $H_P(A_{S,i})$ on the
unmixed pieces and
\begin{equation}
                  H_P(Y_D)+\sum_{v\subseteq D}H_P(J_v)
 \label{e:mixedold}
\end{equation}
on each mixed block. The sum runs over all active internal
nodes in $D$. There the full child observations are used, and
$J_v$ is also the common observation of the raw children.

Write $Z_S$ for the raw words of the $h\ge2$ pieces inside $D$,
and let $E$ be the event that their combined word is constant
with the same bit throughout. By \eqref{e:raw},
\[
 H_P(Y_D)=\sum_{S\subset D}H_P(Z_S)
       -\mathrm{TC}_P((Z_S))-(\log2)P(E),
 \qquad
 \mathrm{TC}_P((Z_S)):=\sum_S H_P(Z_S)-H_P((Z_S)).
\]
For the product law the same identity has zero total correlation
and $b=Q(E)=2^{1-h}\prod_{S\subset D}Q(S\text{ vacant})$.
We show that
\begin{equation}
 \sum_{S\subset D}D(P_{Z_S}\Vert Q_{Z_S})
  +\mathrm{TC}_P((Z_S))-\sum_{v\subseteq D}H_P(J_v)
                  +(\log2)P(E)\ \ge\ (\log2)b .
 \label{e:mixedpay}
\end{equation}

The loss of orientation on $E$ is charged once for the whole block
$D$. We thus need to retain one common observation inside $D$
until this loss has been estimated. Choose a lowest nonfull node
$v_*$ in $D$. Its children are
whole pieces since each has full space and their parent has a proper subspace.
The common observation $J_{v_*}$ is nonzero by the dimension
formula for the two raw restrictions, whose joint determines
the parent. Keep the entire old joint law at this node. In the
tree inside $D$ collapse this node and all other full pieces to
terminal words. Reconstruct upward at every remaining join using
Lemma~\ref{l:join}. This gives a law $P'$ with precisely the
same piece marginals, the same protected joint law, and
$P'(E)\le P(E)$. The total correlation of the pieces is the sum
of the mutual informations at all their binary joins, by entropy
telescoping. All reconstructed joins cost at most their old mutual
information less $H_P(J_v)$. The protected join retains its old cost.
Therefore
\[
  \mathrm{TC}_{P'}((Z_S))-H_P(J_{v_*})\le
  \mathrm{TC}_P((Z_S))-\sum_{v\subseteq D}H_P(J_v).
\]
In particular the left of \eqref{e:mixedpay} dominates
$D(P'\Vert Q_D)-H_{P'}(J_{v_*})+(\log2)P'(E)$ for the product
$Q_D$ of the piece references. The protected common quotient has
two linear realizations, one on each of its full child pieces,
both annihilating constants. A nonzero scalar functional on
the quotient gives a nonzero contrast on each child. Thus the
reference probability of zero on each side is at most $2/3$
by \eqref{e:atom}, since their full observations have positive weight
whenever the mixed cut being considered has positive weight.
The event $E$ implies both realizations are zero. Lemma~\ref{l:tilt}
with $d=\log2$ bounds the last expression below by $3b/4$,
which is greater than or equal to $(\log2)b$. This proves
\eqref{e:mixedpay} simultaneously for the chosen reference
marginals, allowing the auxiliary law $P'$ to depend on $D$.

Finally sum over the weights. For $Q$ the data on distinct
terminals are independent. On mixed $D$ its output entropy is
$\sum_S H_Q(Z_S)-(\log2)b$, and on an unmixed piece it is
$H_Q(A_{S,i})$. The corresponding piece entropy contributions,
before the subtractions $(\log2)b$, sum to $\sum_S G_S(Q_S)$,
by \eqref{e:raw}. The same identity applies to the old marginal
contributions with $P$ instead of $Q$. Lemma~\ref{l:gibbs}
bounds the gain on piece $S$ below by $w_S D(P_{Z_S}\Vert Q_{Z_S})$
because each mixed block sees that piece's full unrooted word.
For $w_S=0$ only nonnegative gain is needed. These gains pay
every divergence in \eqref{e:mixedpay}, exactly once per mixed
block containing its piece. Combining with \eqref{e:mixedold}
proves \eqref{e:comparison}.

\subsection{Integration in time}

For a cut $C_i$ the difference between its full and old observed
ranks at a stage equals $\sum_v j_{v,i}$. Indeed at a node where
the join is not seen the observable is the joint child observable.
Where it is seen the joint unrooted child words acquire one
additional linear coordinate, since both separate child indicators
are in $V_v$. The increment is the same on the full space, and
the common ranks on the full space are zero. The assertion follows
by induction from the cells. The old root observations form an
ordinary test throughout time. For each fixed structural $v$ its
two child kernels pulled back to the restricted $V_v$ are nested
below the split, and their sum is a nested kernel for $J_{v,i}$
containing the constants. Extend it by the whole restricted
space before the split, when that space is just the constants,
and pull back to the original rows. This too is an ordinary test.
Hence the time integral of $H_P(Y_i)$ pays the old observed ranks,
and each integral of $H_P(J_{v,i})$ pays its common ranks.
The integration uses these fixed structural tests, allowing
the pieces in \eqref{e:comparison} to be chosen afresh at each stage. Every chain in \eqref{e:minimax} now has nonnegative
maximal weighted payoff. Lemma~\ref{l:cutmin} and minimax provide
one full weak flag for all tests. This proves
Theorem~\ref{t:restore}.

\section{Large-prime restrictions of close divisors}\label{s:arith}

For an elementary example of a completion, consider the divisors
$15$ and $17$ of $255$. Above a prime cutoff between $5$ and $17$
their logarithmic sums are $0$ and $\log17$, farther apart than
one. Restoring the two smaller primes brings them within
$\log(17/15)$ of each other. For $k$ rows, the smaller primes
can supply a different displacement to every row.

The fixed-range upper argument of Ford, Green and Koukoulopoulos
\cite[Section~4]{fgk} records the first columns that increase the
span and counts the remaining columns in quotient classes. Hou and
Zhao \cite[Lemma~5]{hz} and Mao and Song \cite[Theorem~3.6]{ms}
give probability bounds uniform in a prescribed vector of
translations for the logarithmic random set. Here the lower
prime choices range over all possible completions. We keep them
in the count and mark the height by which the upper restrictions
are already pairwise distinct. This is the condition defining
$s_k$ in Lemma~\ref{l:arithmetic-one-range}.

We use a Poisson random measure $\mathcal P$ with atoms at $\log p$,
for primes $p$, of respective intensities $1/p$. Its counts at
different primes are independent. Copies at an atom are distinguished
when subsets of the process are counted. Write $\mu$ for its intensity
measure. Hou and Zhao use independent geometric prime exponents to
pass between integers and the logarithmic random set
\cite[Section~3]{hz}. We compare the middle range of prime factors
directly with the process on the prime logarithms.

\begin{lemma}\label{l:prime-log-measure}
For $z\ge2$ we have
\[
 A(z):=\sum_{p\le z}\frac{\log p}{p}=\log z+O(1),\qquad
 \sum_{p\le z}\frac1p=\log\log z+B+O(1/\log z)
\]
with a constant $B$. For $v\ge1$, $h\ge0$ and an absolute $C$,
\[
 \frac{h-C}{v+h}\le\mu([v,v+h])\le\frac{h+C}{v}.
\]
The bounds also hold for half-open intervals after increasing $C$.
\end{lemma}
\begin{proof}
For an integer $m\ge1$ the primes in $(m,2m]$ divide
$\binom{2m}{m}\le4^m$. Summing the logarithmic bound on dyadic
intervals gives $\sum_{p\le z}\log p\ll z$. Summing the same bound
over prime powers gives $\sum_{d\le z}\Lambda(d)\ll z$, where
$\Lambda(p^a)=\log p$ for $a\ge1$ and $\Lambda(d)=0$ otherwise.
Since $\sum_{d\mid n}\Lambda(d)=\log n$, summation up to $z$ gives
\[
 z\sum_{d\le z}\frac{\Lambda(d)}{d}
 =\log(\lfloor z\rfloor!)+O(z)=z\log z+O(z)
\]
by integral comparison. The contribution of powers $p^a$, $a\ge2$,
to the sum on the left after division by $z$ is bounded by
$\sum_{n\ge2}\log n/[n(n-1)]$. This proves the estimate for $A$.
Partial summation gives the second estimate, since the integral of
$1/[t(\log t)^2]$ from $z$ to infinity is $1/\log z$. The interval
bounds follow by taking the difference of $A$ at $e^{v+h}$ and $e^v$
and bounding the weight $\log p$ between $v$ and $v+h$.
\end{proof}

\begin{lemma}\label{l:arithmetic-coupling}
Let $N$ be uniform on $1,\ldots,\lfloor x\rfloor$, put
$T=\log\log x$, $G=(\log T)^2$, $y=x^{1/G}$, and
$U=\sqrt{\log T}$. With probability tending to one, the prime
factors of $N$ in $\exp(e^U)<p\le y$ can be coupled exactly to
$\mathcal P$ on the corresponding interval of prime logarithms,
with a single copy of every retained prime. The total number of
omitted prime factors of $N$, counted with multiplicity, is
$o(\log T)$ in probability.
\end{lemma}
\begin{proof}
First couple all valuations at $p\le y$ to independent geometric
variables $Z_p$ with $\PP(Z_p\ge a)=p^{-a}$ for integers $a\ge0$.
Set $q=\prod_{p\le y}(1-1/p)$. Partition these primes into $J=O(T)$
blocks according to the intervals
$(2^{-j-1}\log y,2^{-j}\log y]$ for $\log p$, $j\ge0$, stopping
when the lower end is below $\log2$. In each block the sum of $1/p$
is at most an absolute $C_0$ by Lemma~\ref{l:prime-log-measure}.
Thus its product $q_j=\prod(1-1/p)$ is bounded below by a positive
absolute constant, using $-\log(1-1/p)\le2/p$.

Take an even integer $r=2\lceil\log T\rceil$. For an integer $n$
let $\omega_j(n)$ count the distinct prime divisors in block $j$,
and let $B_j^+$ and $B_j^-$ denote the partial sums of
$\sum_i(-1)^i\binom{\omega_j(n)}i$ through $r$ and $r-1$.
The binomial identity for alternating partial sums gives
$B_j^-\le I_j:=\one_{\omega_j(n)=0}\le B_j^+$, with
$B_j^+-B_j^-=\binom{\omega_j(n)}r$. In particular
\[
 0\le\prod_j B_j^+-\prod_j I_j
 \le\sum_j\binom{\omega_j(n)}r\prod_{h\ne j}B_h^+.
\]
Expanding the first product or any summand on the right in
divisibility indicators gives coefficients of absolute value at
most one on distinct squarefree $d\le D=y^{2r}$. Indeed, block
$j$ contributes at most $r$ primes, each of logarithm at most
$2^{-j}\log y$, so $\log d\le r(\log y)\sum_{j\ge0}2^{-j}=2r\log y$.
Under uniform $n\le X$, $X\ge2$,
the mean of each indicator differs from $1/d$ by $O(1/X)$.
Under independent prime indicators of means $1/p$ its mean is
exactly $1/d$. For these independent variables
$\EE\binom{\omega_j}r\le C_0^r/r!$ and
$q_j\le\EE B_j^+\le q_j+C_0^r/r!$.
The product bound and block independence now give, uniformly in $X$,
\[
 \PP_{n\le X}(p\nmid n\text{ for every }p\le y)
   =q(1+o(1))+O(JD/X).
\]
Indeed $J C_0^r/r!=o(1)$ and the $q_j$ are bounded below.

The independent product $\prod_{p\le y}p^{Z_p}$ takes a given
$y$-smooth value $m$ with probability $q/m$. The exact small-prime
part of $N$ equals $m$ if $m\mid N$ and $N/m$ has no prime factor
at most $y$. Using the last estimate with $X=x/m$, its probability
for $m\le\sqrt x$ is $(q/m)(1+o(1))+O(JD/x)$. The summed absolute
error over these $m$ is $o(1)+O(JD/\sqrt x)=o(1)$ because
$D=x^{2r/G}=x^{o(1)}$. Both small-prime products exceed $\sqrt x$
with probability $O(1/G)$ by Markov's inequality, since each
expected logarithm is at most $\sum_{p\le y}(\log p)/(p-1)\ll\log y$.
The resulting total variation distance tends to zero. Coupling the
common probability masses gives the asserted geometric comparison.

Above $p_0=\exp(e^U)$ the geometric and Poisson laws at $p$ differ
in total variation by $O(p^{-2})$, by comparing the masses at zero
and one and the probability of two or more. The sum of these errors,
including repeated factors in this range, is $o(1)$. Finally
$\PP(p^a\mid N)\le p^{-a}$ and Lemma~\ref{l:prime-log-measure}
give $O(1+U+\log G)$ for the mean number of omitted factors.
Division by $\log T$ and Markov's inequality finish the proof.
\end{proof}

\subsection{An upper bound on the number of continuations}

At $s\ge0$ let $C(s)$ be the largest number of subsets of the
points of $\mathcal P$ in $[1,e^s]$ whose sums lie in one interval
of length one. We may take closed intervals here for an upper bound.

\begin{lemma}\label{l:arithmetic-one-range}
Fix $k\ge2$ and $s_k<\rho<1$. With probability
$1-O_{k,\rho}(e^{-c s})$, for some $c=c(k,\rho)>0$ and $s\ge1$,
there is no family of $k$ subsets of $\mathcal P\cap[1,e^s]$
whose sums have diameter at most one and whose restrictions to
$(e^{\rho s},e^s]$ are all distinct. In particular
$C(s)\le(k-1)C(\rho s)$ on this event.
\end{lemma}
\begin{proof}
For a fixed small $\xi>0$, except with probability
$O_\xi(e^{-c_\xi s})$, every interval of log coordinates
$(a,b]\subseteq[0,s]$ contains at most $b-a+\xi s$ points.
The same bound with an extra absolute constant covers the whole
range. Indeed the mean of each prefix count from $1$ to $e^b$ is
$b+O(1)$ by Lemma~\ref{l:prime-log-measure}. For its mean $m(b)$
and count $N(b)$, $\exp(\lambda N(b)-(e^\lambda-1)m(b))$
is a mean-one martingale by independent increments. Applying its
maximal inequality for both signs of a sufficiently small fixed
$\lambda$, depending on $\xi$, bounds the largest centered prefix
deviation by $\xi s/3$ with the claimed exceptional probability.
Subtract prefixes to obtain the assertion for large $s$.

For each point of a witness write its $k$ Boolean membership bits
as a column. Read the points in decreasing order, breaking ties in
any order, and mark a pivot whenever its column enlarges the span
of the constant and all previous columns. Let $u_1,\ldots,u_r$ be
the marked sizes, where $r\le k-1$, and $V_j$ the span through
pivot $j$, with $V_0=\Q\one$. Round $(\log u_j)/s$ upwards to the
grid of step $1/\lceil s\rceil$ ending at $1$, obtaining $c_j$,
and set $c_{r+1}=0$. Each residual point in
$(e^{s c_{j+1}},e^{s c_j}]$ has column in $V_j$, and columns above
the first layer are constant. There is no prime logarithm equal to
$1$. The span of the columns with point size exceeding $e^{\rho s}$
already distinguishes all rows, so some separating $V_j$ has
$c_j\ge\rho$. Let $P_j$ be the residual empirical law in layer
$j$, using the zero word for an empty layer.

On the count event there are $s^{O_k(1)}$ possibilities for the
pivot columns, their rounded heights, and the counts of individual
Boolean words in every layer (for bounded $s$ one may replace $s$
by $2+s$). Uniformly for these data the minimum over tests of the
zero-endpoint payoff \eqref{e:pay} is at most $-\delta$ for some
$\delta=\delta(k,\rho)>0$. To see this, fix one of the finitely
many ambient flags. Its ordered clocks and stage probability
simplexes are compact, including the closed condition that a
separating space be born at or above $\rho$. For tests it suffices
to keep one representative of each attainable list of dimensions
and partitions of the finite stage alphabets into cosets. The
minimum of their continuous payoffs is continuous. A nonnegative
maximum on this compact set would give $s_k\ge\rho$.

Fix a type and its minimizing test $W$. Integrate the pivot
positions against $\prod_{i=1}^r d\mu(u_i)$, leaving an independent
residual Poisson process. This insertion identity bounds the
probability of a witness by the expected number of ordered tuples
of distinct point copies that can serve as pivots. To verify the
identity, condition on a Poisson count $n$ and use the $n$ independent
positions of law $\mu/\mu([1,e^s])$. Multiplying by the falling
factorial $(n)_r$ and summing replaces $n$ by $n-r$ in the Poisson
series and gives the displayed product measure. This argument also
holds at atoms. Only residual configurations satisfying the upper
count bounds need be considered, since deletion decreases counts.

Given such residual positions, specifying the coset at every
position in layer $j$ has at most
$\exp(|S_j|H(P_j\bmod W_j))$ possibilities by the multinomial
bound, where $S_j$ is the set of residual point copies in the layer.
Thus the logarithm of the total coset cost is at most
$s\sum_j(c_j-c_{j+1})H(P_j\bmod W_j)+O_k(\xi s+1)$.
Subtract a fixed Boolean representative for each coset.
Modulo constants, the remaining residual sum has components in
a basis adapted to the nested $W_j$. A coordinate first appearing
in $W_j$ is bounded in absolute value by $O_k(s e^{s c_j})$
because it receives contributions only from that layer and the lower ones.
The bases are fixed for the finite list of test representatives.
Covering these ranges with unit cubes consequently costs at most
\[
 \exp\!\left(s\sum_j c_j(\dim W_j-\dim W_{j-1})
                    +O_k(\log(2+s))\right).
\]
This is a cover of real sums and uses no lattice condition on the
prime logarithms.

In one such cube the residual sums modulo constants have bounded
uncertainty in each standard row contrast against the first row.
Some $r$ such contrasts of the pivot columns give an invertible
integer matrix. Bounded diameter of the final sums then places
each pivot $u_j$ in an interval of length $O_k(1)$, with location
determined by the coset list, the cube and the residual positions.
Also $u_j\ge e^{s c_j-1}$ by rounding. By
Lemma~\ref{l:prime-log-measure} the product of the $r$ pivot
intensities on these intervals is $O_k(\exp(-s\sum_jc_j))$,
including bounded ranges at the bottom of the grid.
Since each pivot raises dimension by one, telescoping gives
\[
 \sum_j c_j(\dim W_j-\dim W_{j-1})-\sum_j c_j
  =-\sum_j(c_j-c_{j+1})\dim(V_j/W_j).
\]
Summing over the types and integrating the residual process yields
$\exp(-\delta s+O_k(\xi s+\log(2+s)))$ on the count event.
Choosing $\xi$ small enough proves the probability estimate.
Finally each fixed upper restriction of a family in one interval
allows at most $C(\rho s)$ lower completions by translation.
\end{proof}

\subsection{Iteration and the normal order}

We now iterate the one-range estimate. Use the coupling in Lemma~\ref{l:arithmetic-coupling} and set
$S=\log\log y=T-\log G$. A divisor of the omitted part of $N$ can
be fixed in at most $2^{o(\log T)}$ ways in probability. After it
has been fixed, the retained divisors are subsets of the coupled
point process and their logarithms are the corresponding subset sums.
Thus $\log\Delta(N)\le\log C(S)+o(\log T)$ in probability.

For fixed $s_k<\rho<1$ apply Lemma~\ref{l:arithmetic-one-range}
at the successive times $S,\rho S,\rho^2 S,\ldots$ down to the
last time at least $U$. There are at most
$1+\log(S/U)/\log(1/\rho)$ such times. The total exceptional
probability is $O_{k,\rho}((1+\log T)e^{-cU})=o(1)$. The number
of points up to $e^U$ is $O(U)$ with probability tending to one,
by the Poisson exponential bound. Consequently
\[
 \frac{\log\Delta(N)}{\log T}
 \le\frac{\log(k-1)}{\log(1/\rho)}+o(1)
 \qquad\text{in probability as an upper bound.}
\]
Let $\rho\downarrow s_k$ and take the infimum over $k$ with $s_k<1$.

For the other direction fix a strict separating flag on $k$ rows
at endpoint $b>0$ and choose $0<b'<b$. We may first mix with
an arbitrarily small proportion of the auxiliary strict clock, so
that all birth times are distinct and the endpoint still exceeds
$b'$. By
Lemma~\ref{l:realization}, in $(e^{b's},e^s]$ the prime-log
process has $k$ distinct subsets of sums of diameter at most a
fixed constant $L$ with probability tending to one as $s\to\infty$.
Here $L$ may depend on the fixed flag and on $b'$. Use the disjoint
intervals with upper times $s_j=S(b')^j$, $j=0,\ldots,m-1$,
where $m=\lfloor\log(S/U)/\log(1/b')\rfloor$. Every lower time
is at least $U$ and the least upper time tends to infinity. Hence
the expected fraction of failed intervals tends to zero, by
uniformity of an ordinary limit over all $s\ge U$. Markov's
inequality shows that the number of successes is $m-o(m)$ in
probability. On the coupling event all their primes divide $N$
with exponent one. Taking one subset from each successful interval
and uniting them gives $k^{m-o(m)}$ distinct divisors by disjoint
supports and unique factorization. Their logarithms have diameter
at most $Lm$. This set is covered by at most $\lceil Lm\rceil+2$
intervals of the form $(u,u+1]$, so some unit interval contains
at least $k^{m-o(m)}/(\lceil Lm\rceil+2)$ of these divisors.
Since $\log(S/U)=(1+o(1))\log T$, this gives the lower bound
$\log k/\log(1/b')$ for $\log\Delta(N)/\log T$ in probability.
Take $b'$ up to the endpoint and then take the supremum over the
strict flags.

Combining the two arithmetic estimates with the
comparison in Section~\ref{s:open} proves \eqref{e:alpha}.
With probability $1-o(1)$ we have $N>\sqrt{x}$, so that
$\log\log N=T+O(1)$ and $\log\log\log N=\log T+o(1)$.
This proves the convergence in density in Theorem~\ref{t:main}.
For the equivalent pointwise statement let $E_j$ be the exceptional
set at tolerance $1/j$. Choose increasing integers $x_j$ such that
$\#(E_j\cap[1,x])\le x/j$ for every $x\ge x_j$. Delete $E_j$
on $[x_j,x_{j+1})$. For $x_j\le x<x_{j+1}$ all the deleted
integers up to $x$ belong to $E_j$, since $E_i\subseteq E_j$
for $i\le j$. The deleted set has density zero and the exponent
converges on its complement.

\section{Effective evaluation of the variational exponent}\label{s:eval}

The supremum in \eqref{e:alpha} provides approximations from
below and the infimum provides approximations from above. To turn
these formulas into an algorithm, two errors must be controlled.
For a fixed $k$, a strict flag certifies feasibility and compactness
gives a uniform certificate of failure beyond the endpoint.
For varying $k$, Lemma~\ref{l:opening} bounds the loss of using
only $s_k$. We give the two estimates in that order.

For a rational trial $d>b_k$, the minimum payoff over tests is
strictly negative on every separating flag with endpoint $d$.
A nonnegative minimum would give a weak flag, hence strict
endpoints approaching $d$. The ordered clocks at fixed endpoint
and stage probability simplexes form compact rational polytopes,
so there is a uniform negative gap. Subdivide them into rational
boxes, discarding boxes disjoint from the domain by rational
linear feasibility. Enclose the finitely many payoffs by rational
intervals with outward rounding. The oscillation of $-x\log x$
on an interval of width $h\le1$ is at most $h(2-\log h)$. For
logarithms of positive rationals one can use range reduction
and the convergent series for $\log((1+z)/(1-z))$. Refining the
boxes and the arithmetic eventually certifies negativity everywhere.
For $d<b_k$, enumeration of rational strict flags with endpoint
larger than $d$ eventually succeeds, by continuity and the finite
list of proper tests.

For $s_k$, impose endpoint zero and require the first separating
space to have birth at or above the rational trial $d$. These
are finitely many compact rational domains, establishing also
attainment in its definition. The same negative-gap search succeeds
for $d>s_k$. For $d<s_k$, mix a weak flag separating strictly
above $d$ with the auxiliary strict clock. It still separates
strictly above $d$ and has a positive endpoint. Rational strict
flags of this kind can be enumerated, and each gives a valid
lower certificate by filling to zero. Running both strict
comparisons on all rational trials in parallel gives brackets
of arbitrarily small positive width for each finite endpoint, even
if it coincides with one of the rational trial values.

For an explicit stopping rule, restriction gives a flag with
exactly $k$ rows distinguished at $s_k$. The proof of
Lemma~\ref{l:opening} applies to those rows, yielding
$b_{\lfloor k/5\rfloor}\ge s_k/8$ for $k\ge10$.
The variational formula \eqref{e:alpha} gives \eqref{e:singlek}. Write
$D=\log(1/s_k)\ge\log(k/2)/a_0$, by the three-row depth bound.
The width in \eqref{e:singlek} is at most
\begin{equation}
 \frac{\log10}{D}+\frac{(\log k)\log8}{D^2}
 \ \le\
 \frac{a_0\log10}{\log(k/2)}
 +\frac{a_0^2(\log k)\log8}{\log^2(k/2)} .
 \label{e:width}
\end{equation}
Choose $k$ large enough that the last expression is less than
$\varepsilon/2$, with rational interval arithmetic. For this one
$k$, compute successively finer certified brackets for $s_k$.
Use the lower endpoint in the left side of \eqref{e:singlek}
and the upper endpoint in its right side, rounding both ratios
outwards. Since $0<s_k<1$ and their actual separation is already
less than $\varepsilon/2$, the resulting rational bounds reach
width less than $\varepsilon$ in finite time. Theorem~\ref{t:restore}
also allows the finite optimization for $s_k$ to be restricted to
full partition clocks.

\section{A certified lower bound from sixteen types of splits}
\label{nl:section}

In the binary flags of Ford, Green and Koukoulopoulos every cell
at a given level divides at the same height
\cite[Definition~3.8 and Sections~9--10]{fgk}. The lower formula
in \eqref{e:alpha} also allows the two branches of a cell to
divide at different heights. The problem is to produce many leaves
above a small height $\theta$ while every contrast still earns
enough entropy during its lifetime. Two siblings born at different heights face different
opportunities for gaining entropy. When the earlier one is born,
its sibling has yet to split. When the later one is born, its
sibling has already split. To choose the relative heights of a
node's children, we remember whether the node was the earlier or
later sibling and the preceding three turns on its path. This gives
eight sibling pairs, or sixteen types, with the same rules recurring
down the tree.
Correlating the choices within a pair can increase the entropy at
its parent. One entropy budget per type certifies feasibility,
and a spectral-radius equation counts the resulting leaves.

\begin{proposition}\label{nl:result}
Let $d$ be the positive solution of $\rho(M(d))=1$, where the
eight-by-eight matrix $M$ is given in \eqref{nl:matrix}. Then
\[
      0.35553850187<d<0.35553850189,\qquad d\le\alpha .
\]
\end{proposition}

\subsection{The split tree and its parameters}

A node denotes a split. Its type is $i\in\{0,\ldots,15\}$ and its
height is $\tau$. The root is of type $0$ at height $1$. The left
and right child splits have respective types and heights
\[
 L(i)=2(i\bmod 8),\quad R(i)=L(i)+1,\qquad
 A_i\tau,\quad B_i\tau,\qquad B_i=A_i\kappa_{i\bmod8}.
\]
Write $K_{2j}=\kappa_j$, $K_{2j+1}=A_{2j}/\kappa_j$ for $0\le j<8$.
At a nonroot node of type $i$, $K_i$ is the height of the next
split in its sibling's subtree, in units of its own split height. For an even type it is
the birth of the younger sibling. For an odd type it is the first
child split of the older sibling, since $\max_i A_i<\min_j\kappa_j$.
The baseline probability used by an active type $i$ at local time
$s=t/\tau$ is
\[
 P_i(s)=q_{i,m_i(s)},\qquad
 m_i(s)=\one_{s<K_i}+\one_{s<A_i}+\one_{s<B_i}.
\]
The adjustments $h_i$ and the baseline parameters $q_{im}$ are
extended with period eight in $i$.
All displayed parameters are exact rationals.
\[
\begin{array}{c|rrrr}
j&A_{2j}&A_{2j+1}&\kappa_j&\delta_j\\ \hline
0&.17008802076989&.17445386687652&.68401572231443&.003668869499\\
1&.16936398844577&.17388142632538&.66471379404345&.004076944109\\
2&.17009530835310&.17438951025087&.68921875368253&.003631054748\\
3&.16945123639388&.17388691615311&.66892994497063&.004049115897\\
4&.17008831268608&.17445384071325&.68393223772060&.003668869499\\
5&.16936298668175&.17388152290720&.66497954537033&.004076944109\\
6&.17009683577282&.17438937491771&.68877984539362&.003631054748\\
7&.16945134069206&.17388690617701&.66890195988505&.004049115897
\end{array}
\]
{\small
\[
\begin{array}{c|rrrr|r}
i&q_{i0}&q_{i1}&q_{i2}&q_{i3}&h_i\\ \hline
0&.367906599960&.378382774055&.378609457829&.382868703940&.005753434443\\
1&.367906988749&.370803262020&.371785782457&.371946413457&.004550810468\\
2&.367995493671&.379097176841&.379119849971&.383989643006&.006430082793\\
3&.367881403353&.371076969007&.372210929174&.372308612632&.004990980423\\
4&.367879728139&.378327161781&.378614497836&.382778568605&.005690650612\\
5&.368044463559&.370926681912&.371881452912&.372061709237&.004521571727\\
6&.367880216220&.378897274167&.378922330993&.383794414115&.006384472359\\
7&.367881830045&.371054243031&.372198860777&.372294368723&.004974131778
\end{array}
\]}

\subsection{A law on the active cells}

At each $t>0$ only finitely many splits have height at least $t$.
A pattern is the unrooted Boolean word on the current cells of a
subtree. A split is eligible if it is active and its two child
patterns are vacant. An eligible split can choose a vacant parent
or one of two opposite orientations. If exactly one child pattern
is vacant, its absolute bit is fair. If neither is vacant the two
raw representatives determine the joined word.

Generate a pair of sibling patterns from the bottom up. Generate
the proper descendant forests of the two siblings independently,
using this same pair rule within each forest. Let $l,r$ be the
eligibility probabilities of the siblings and $p,q$ their baseline
probabilities at their respective local times. Given eligibility
of the left sibling, its conditional probability of choosing a
vacant rather than an oriented parent is
\[
 p_1=p-h_L(1-r)\quad\hbox{if the right is eligible},
 \qquad p_0=p+h_Lr\quad\hbox{otherwise},
\]
where $L$ is the left sibling's type. On the right, of type $R=L+1$,
use $q_1=q-h_R(1-l)$ and $q_0=q+h_R l$.
At double eligibility use the product table for the two vacancy
indicators with marginals $p_1,q_1$, subtract $\delta_{L/2}$ from
both diagonal entries and add it to both off-diagonal entries.
Conditional on the indicators, all nonzero orientations are fair
and independent. In every other case use conditionally independent
joins. An inactive split supplies a vacant pattern and is ineligible.
At the root use its own baseline probability directly, with a fair
absolute bit on vacancy. The parameters give adjusted probabilities
in $[.35,.40]$, and $\max_j\delta_j<.0041$, so every entry of each
modified table is positive.

Let $v_i(s),z_i(s)$ be the vacancy and eligibility probabilities of
a split of type $i$. For $s>1$ these are $(1,0)$. For $s\le1$ write
\[
 (x,l)=(v_L,z_L)(s/A_i),\qquad(y,r)=(v_R,z_R)(s/B_i),
 \qquad L=L(i),\ R=R(i).
\]
Use $p=P_L(s/A_i)$, $q=P_R(s/B_i)$ in the pair rule above. Averaging
over the independent opposing forest gives $(1-r)p_0+rp_1=p$ and
the analogous identity for $q$. Hence
\begin{equation}
 z_i(s)=xy-lr\{pq-p_1q_1+\delta_{i\bmod8}\},\qquad
 v_i(s)=P_i(s)z_i(s).
 \label{nl:vacancy}
\end{equation}
If both children are active their simultaneous vacancy requires
double eligibility and has probability $lr(p_1q_1-\delta_{i\bmod8})$.
If a child is inactive, $lr=0$ and simultaneous vacancy has
probability $xy$. These cases prove the recursion.

At time $t=1/10$ the two children of the root have split, and
their four children have yet to split. Both sibling choices are
eligible. Relative to independent choices with the same marginals,
the adjustment $\delta=\delta_0$ moves mass $\delta$ from double
vacancy and $\delta$ from double nonvacancy to the two cases
with exactly one vacancy. Thus the root gains at least
$\delta(2\log2-\log3)=\delta\log(4/3)$ in conditional entropy,
since a single vacancy gives one fair bit and a double vacancy
gives at most three outcomes. The entropy loss in combining the
children is their mutual information, at most $\delta^2/(.35\cdot.60)^2<\delta/10$ by
$\log x\le x-1$.
The adjustment therefore increases the entropy of the unrooted
word on these four cells.

\subsection{Entropy allocated to a split}

Use the leftmost descending path of a cell as its representative.
The contrast $Z_v$ at split $v$ is the bit of its right child's
representative minus that of its left child's representative.
Order the active splits by decreasing height, breaking ties
consistently. Given the representative bit at the root, each
contrast determines the other child representative bit. The
contrasts therefore form linear coordinates modulo constants. Put
\[
 G(p)=-p\log p-(1-p)\log((1-p)/2).
\]
Write $I_j(p,q)$ for the mutual information of the two indicators
with marginals $p,q$ and the $\delta_j$ adjustment. For split type
$i$, let $i'$ be its sibling type, $j=\lfloor i/2\rfloor$, and set
$s'=s/\kappa_j$ if $i$ is even, $s'=s\kappa_j$ if it is odd. Write
\[
 z=z_i(s),\quad z'=z_{i'}(s'),\quad
 u_0=P_i(s)+h_i z',\quad
 u_1=P_i(s)-h_i(1-z'),\quad
 u'_1=P_{i'}(s')-h_{i'}(1-z).
\]
The $x,y$ below are the vacancy probabilities of the two children
of the target split in \eqref{nl:vacancy}. Its contrast entropy
conditional on all later active contrasts is at least
\begin{align}
 H_i(s)={}&z\{(1-z')G(u_0)+z'G(u_1)\}
                   +(x+y-2z)\log2  \notag\\
          &-\one_{\{i\ {\rm even},\ s<\kappa_j\}}
                   zz'I_j(u_1,u'_1).
 \label{nl:entropy}
\end{align}
Values at the finitely or countably many boundaries never affect
our integrals.

To see which observations enter this conditional entropy, every
strict ancestor of the target is earlier, as is every sibling of
such an ancestor. Indeed a descendant split lies at height at most
$\max_i A_i$ times its ancestor's height, and $\max_i A_i<\min_j
\kappa_j$. At each ancestral pair both joins are therefore absent
from the conditioning. After marginalizing these joins the two
proper descendant forests are independent by construction. Descending
the ancestral path shows that outside observations do not change the
conditional entropy within the target pair.

Now condition additionally on the two complete proper descendant
forests of this pair. They determine the patterns below each
sibling. On target eligibility the conditional entropy of its join
is the first term of \eqref{nl:entropy} after averaging, since
eligibility events in the two forests are independent. Exactly one
vacant child contributes a fair bit, giving the second term. A
younger active sibling has its own contrast among the conditioning
variables. Given the forests, the resulting debit is supported on
double eligibility. There the zero indicators of the two sibling contrasts have mutual
information $I_j$, and the conditional fair signs are independent.
Only the older (even) sibling incurs this debit. At the actual root
the same bound, using a virtual sibling, follows from
$G(P_i(s))\ge(1-z')G(u_0)+z'G(u_1)$ and nonnegativity of the debit.

\subsection{The rational entropy certificate}

Set $c_0=1/4000$. The following table gives a lower certificate for
$\int_{c_0}^1 H_i(s)\,ds-1$. Its last row gives the floors of the
certified lower estimates multiplied by $10^{12}$.
\[
\begin{array}{c|rrrrrrrr}
i&0,8&1,9&2,10&3,11&4,12&5,13&6,14&7,15\\ \hline
\text{intervals}&56&60&56&60&56&60&56&60\\
\text{scaled surplus}&29&29&29&30&30&30&29&29
\end{array}
\]
Here are precise instructions for verifying the table by rational
arithmetic. Start the type-$i$ tree at height one and the sibling
tree at relative height $\kappa_j$ (even $i$) or $1/\kappa_j$ (odd
$i$). At a node of type $a$ and height $w\ge c_0$, record $w,wK_a$
and recur at $wA_a,wB_a$. Stop if $w<c_0$. Restrict the union of
the recorded points to $[c_0,1]$ and adjoin both endpoints. These
points include all discontinuities in \eqref{nl:vacancy} and
\eqref{nl:entropy}, since both child births are recorded and
$K_a,A_a,B_a<1$. Evaluate at the rational midpoint of each interval
and sum its length times the lower bound for \eqref{nl:entropy}.

For positive rational $y$ put $v=(y-1)/(y+1)$. The $N$-term
sum $2\sum_{m=0}^{N-1}v^{2m+1}/(2m+1)$ approximates $\log y$ with
absolute error at most
$2|v|^{2N+1}/((2N+1)(1-v^2))$. Use $N=40$ for each logarithm
in $G$, $I_j$ and $\log2$, rounding the argument first and the
result last outwards on the mesh $10^{-22}$. For the lower bound
on $G$ use upper bounds on $\log p$ and $\log((1-p)/2)$.
For the upper bound on $I_j$ sum each positive cell mass times
the upper bound on the log of its ratio to its product mass.
Use the lower bound on $\log2$ in \eqref{nl:entropy}. These are the
operations of the accompanying \texttt{sixteen\_certificate.py}.

The table also controls termination at any intermediate height.
Let $\delta_* =\max_j\delta_j$. The inequality $\log x\le x-1$
gives
\[
 I_j(p,q)\le \frac{\delta_*^2}{p(1-p)q(1-q)}
          \le \frac{\delta_*^2}{(.35\cdot .60)^2}
 \quad(.35\le p,q\le .40).
\]
Using the same series bounds gives
\[
 \min\{G(.35),G(.40)\}-\frac{\delta_*^2}{(.35\cdot .60)^2}>1,
 \qquad (1+\max_{a,m}q_{am})\log2<1.
\]
Concavity of $G$ thus gives $H_i(s)\ge0$ for every $0<s<1$.
When $A_i<s<1$ both children are unsplit, so $z=1$ and $H_i(s)>1$.
When $s<A_i$ the left child is active, whence $x\le\max q_{am}$
and $H_i(s)\le(x+y)\log2<1$ by $G\le\log3<2\log2$.
Therefore $\int_r^1(H_i(s)-1)\,ds$ is positive for $A_i\le r<1$,
and on $[c_0,A_i]$ its minimum is at $c_0$. Together with the
table and $H_i\ge0$ this proves
\begin{equation}
 \int_r^1 H_i(s)\,ds>1-r\quad(0\le r<1),\qquad
 \int_r^1 H_i(s)\,ds>1\quad(0\le r\le c_0).
 \label{nl:budget}
\end{equation}

\subsection{Weak flags from finite truncations}

Fix $0<\theta<1$ and perform all splits of height at least
$\theta$. Take the leftmost path in each of the $k$ frontier cells
as a labelled row. At time $t\in[e,1]$, with $e=c_0\theta$, first
draw the active tree law above and read it on these rows. This
gives a Boolean word in the full partition space on the current
target cells. A target contrast equals the corresponding contrast
in the full active tree by the choice of paths. Conditioning just
on the later target contrasts can only increase the lower bound
\eqref{nl:entropy}. Average the raw laws within each stage of the
target flag. Entropy concavity preserves every integrated lower
estimate for the quotient entropy of a fixed test on that stage.

Here is the bookkeeping for an arbitrary nested test $W$, modulo
constants. Use the chronological contrast coordinates and take the
rightmost nonzero coordinate as the pivot in elimination. The set
of pivot indices consists of all rightmost nonzero indices of
vectors in $W$, and increases when $W$ increases. The quotient has
one coordinate per nonpivot index $i$, equal to $Z_i$ plus a linear
combination of later contrasts. By the reverse entropy chain rule
its entropy is at least the sum over nonpivots of $H(Z_i\mid
\text{all later contrasts})$. Indeed conditioning on the later
contrasts determines the later quotient coordinates and the additive
term in the $i$th one.

A nonpivot index born at height $\tau$ and first pivoted at
$r\le\tau$ costs rank $\tau-r$ up to that closure. It earns at
least $\tau\int_{r/\tau}^1H_i(s)\,ds$, covering the cost by
\eqref{nl:budget}. An index still unpivoted at $e$ costs
$(\tau-e)+e=\tau$ including the terminal debit, and $e/\tau\le
c_0$ gives the stronger estimate in \eqref{nl:budget}. Immediate
pivots cost zero, including at tied births. Summation proves
weakness at the positive endpoint $e$. Strict approximation on this
finite flag, as in Section~\ref{s:flags}, therefore gives $b_k\ge c_0\theta$.

\subsection{Growth of the tree}

For $0\le p<8$ put $j=2p\bmod8$. The only nonzero entries in
row $p$ of $M(d)$ are
\begin{equation}
        M_{p,j}(d)=A_{2p}^d,\qquad
        M_{p,j+1}(d)=(\kappa_p A_{2p+1})^d .
 \label{nl:matrix}
\end{equation}
Its edges append a binary digit to the three-digit state $p$ and
discard the first digit. Thus $M^3$ is positive. Its spectral
radius is continuous and strictly decreasing from two at $d=0$
to zero as $d\to\infty$, by positivity and the strict decrease of
the edge weights. At the unique root take $D>0$ with $M(d)D=D$,
and put $w_i=A_i^d D_{i\bmod8}$. Then
\[
 w_{2p}+\kappa_p^d w_{2p+1}=D_p,\qquad
 A_i^d w_{L(i)}+B_i^d w_{R(i)}=w_i.
\]
Consequently the mass $\tau^d w_i$ of each split is conserved
by its replacement with its two child splits. Every frontier cell
is identified by its next split, of height less than $\theta$.
Conservation gives $k\ge w_0/(\theta^d\max_i w_i)$. Applying
the lower formula in \eqref{e:alpha} to $k$ with $b_k\ge c_0\theta$ and letting
$\theta\downarrow0$ proves $\alpha\ge d$.

To enclose the root, take the rational vector
\[
\begin{split}
 D_*=(\,&1.00000000000000,\ .99507721338637,\ 1.00122719927539,\ .99616212670301,\\
       &.99998002103143,\ .99514192367964,\ 1.00112266506735,\ .99615533827001).
\end{split}
\]
At $d_-=.35553850187$ the least coordinate of $M(d_-)D_*-D_*$
exceeds $2.08\cdot10^{-11}$. At $d_+=.35553850189$ the greatest
coordinate of $M(d_+)D_*-D_*$ is less than $-1.78\cdot10^{-11}$.
For these inequalities compute each $x^d=\exp(-d\log(1/x))$
using the same outward logarithm bounds with $N=120$. The
arguments $u=d\log(1/x)$ lie in $[0,2]$, so the Taylor sums for
$\exp(-u)$ through orders $35$ and $34$ bound it below and above.
The componentwise inequalities, tested against a positive left
Perron eigenvector, put $\rho(M(d_-))>1>\rho(M(d_+))$. This proves
Proposition~\ref{nl:result}.

\section{A certified upper bound for the exponent}
\label{nu:section}

Put $a=35563/100000$, $C=1499091/10$, and $\ell=\log 2$.
We prove
\begin{equation}
 \alpha\le a_*:=\frac{a\ell}{\ell-\log(1-1/C)}
                  <0.3556265775100 .                         \label{nu:bound}
\end{equation}
The lower bound counted the leaves of one feasible tree. To bound
all such trees, assign each split of height $T$ the weight $T^a$.
A bound $C$ for the sum of these weights implies that at most
$C c^{-a}$ splits occur at heights $T\ge c$. There is an additional
gain because the same inequality holds inside every subtree after
rescaling. This accounts for the improvement from $a$ to $a_*$
in \eqref{nu:bound}. The entropy calculation below establishes the
uniform estimate on the sum of weights.

\subsection{Entropy on the active tree}

By Theorem~\ref{t:restore} it suffices to work on a weak full
partition flag. Write $T_b$ for the split height of an internal
vertex $b$, assigning height zero to terminal vertices. Only
positive heights matter. A split is active at time $t<T_b$.
For each $b$ observe the unrooted word $Y_b$ on its rows, using
the zero observation before its split. At an internal vertex order
its children $L,R$ so that
\[
 u=T_L/T_b\ge v=T_R/T_b,\qquad U=u^{1-a},\quad V=v^{1-a}.
\]
Set $w_b=T_b^{a-1}$ and $S_b=(T_b/T_{\rm parent})^{1-a}$,
taking $S_b=0$ at the root. The three-row estimate in the proof
of finiteness of $M(c)$ gives $u,v\le h=0.65$, since its exact
ratio $(\log3-1)/(\log3-1+2-\log7)$ is less than $h$.
Also $S_b\le0.76$, since $14(1-a)>9$ and $h^9<0.76^{14}$.
The choice $w_b=T_b^{a-1}$ gives a unit rank cost over the lifetime
of $b$ total weight $w_bT_b=T_b^a$. This lets us sum local entropy
inequalities to control the number of splits. Every terminating
decimal below is an exact rational number.

For a linear observation $Y$ on the current full space let $e_Y$
denote the integral of its entropy minus its rank. Weakness gives
$e_{Y_b}\ge0$ and $e_{(Y_L,Y_R)}\ge0$. The full space makes the
two child images a direct product. The rank on the parent exceeds
their joint rank by one while the parent is active. Thus
\[
 e_{Y_b}-e_{Y_L}-e_{Y_R}
  =\int_0^{T_b}\big[H(Y_b\mid Y_L,Y_R)-1-I(Y_L;Y_R)\big]dt,
 \qquad
 e_{(Y_L,Y_R)}=e_{Y_L}+e_{Y_R}-\int I(Y_L;Y_R)dt .
\]
Give an internal vertex the nonnegative weight
$\beta_b=w_b-w_{\rm parent}$, with parent weight zero at the
root, and assign terminal vertices $\beta_b=0$. In the sum
$\sum_b\beta_b e_{Y_b}$ replace $m_b=\min(\beta_L,\beta_R)$
copies of the separate child tests by their joint test. Each child
is used at its unique parent. The resulting sum is nonnegative.
Telescoping it using the displayed identities gives
\begin{equation}
 0\le\sum_b w_b\int_0^{T_b}
   [H(Y_b\mid Y_L,Y_R)-1-(1+m_b/w_b)I(Y_L;Y_R)]dt .       \label{nu:entropy}
\end{equation}
The coefficient of information at two active children is $1/U$.
With at most one active child their mutual information is zero.

Write $p_b=\PP(Y_b=0)$, $\theta_b=S_b\ell$, $E(y)=\log(2+e^y)$,
$q(y)=e^y/(2+e^y)$, and $D(y)=2\ell-E(y)$. Add
$\theta_bp_b-\ell\sum_{i\ {\rm active\ child}}p_i$ to the
integrand at each active vertex. These terms cancel in
\eqref{nu:entropy} since $w_i\theta_i=w_b\ell$.
For $0\le r\le1$, $s>0$, define
\[
 A_{s,r}(d)=-s^{-1}\log(1-r+re^{-sd}),\qquad
 A_{0,r}(d)=rd,\qquad A_s(d)=A_{s,q(s\ell)}(d).
\]
Let $K_-(p,r)=p\log(p/r)+(1-p)\log((1-p)/(1-r))$ for $p<r$
and let $K_-(p,r)=0$ for $p\ge r$, with limiting conventions.

At a fixed time propagate numbers $r_b$ upwards through the active
tree. Suppose the active children already supply deductions
$K_-(p_i,r_i)$ in their own weight units. For $y\le\theta_b$ put
\[
 \Gamma_b(y)=
 \begin{cases}
 E(y),&\text{no active children},\\
 \ell-A_{U,r_L}(D(y)),&\text{one active child},\\
 -A_{V,r_R}(A_{U,r_L}(D(y))),&\text{two active children}.
 \end{cases}
\]
After applying the deductions of the children, the local
contribution at price $\theta_b$, in units of $w_b$, is bounded by
\begin{equation}
  \Gamma_b(\theta_b)-1-K_-(p_b,r_b),
  \qquad r_b=\Gamma_b'(\theta_b).                           \label{nu:reserve}
\end{equation}
To prove this, first use price $y$ and retain the universal $-1$
separately. Two vacant children allow three unrooted extensions.
Their conditional entropy plus $y$ times probability of parent
vacancy is at most $E(y)$ by exponential tilting. A single vacant
child allows two extensions and two nonvacant children determine
the parent. With one active child of vacancy probability $x$,
these bounds, including the child price, give $\ell-D(y)x$.
With two active children they give $-D(y)\PP(B_L B_R=1)$, where
$B_i=\one_{Y_i=0}$. Here $0\le D(y)\le\ell$. Data processing
gives $I(Y_L;Y_R)\ge I(B_L;B_R)$. For any marginal $p_i>r_i$,
independently thin its indicator to marginal $r_i$. This only
reduces the intersection and the mutual information, and turns
its zero one-sided cost into a zero ordinary divergence.
Thus all indicator costs can now be taken to be ordinary
Bernoulli divergences. The minimum of $dx+s^{-1}D(\mathrm{Ber}(x)
\Vert\mathrm{Ber}(r))$ is $A_{s,r}(d)$. For two children combine
$I(B_L;B_R)$ and the left marginal divergence into the conditional
divergence of $B_L$ given $B_R$, and minimize first conditionally
on $B_R$, then over its law. The coefficients are respectively
$1/U$ and $1/V$. This proves the upper bound $\Gamma_b(y)-1$
at every price $y\le\theta_b$.

Now $A'_{s,r}(d)=re^{-sd}/(1-r+re^{-sd})$ and its derivative
is $-sA'_{s,r}(1-A'_{s,r})$. With two active children let
$z=q(y)$, $l=A'_{U,r_L}(D(y))$ and $m=A'_{V,r_R}(A_{U,r_L}(D(y)))$.
Then $Q(y)=\Gamma_b'(y)=zlm$ and
\[
 Q'(y)/Q(y)=(1-z)+Uz(1-l)+Vzl(1-m)\le1-Q(y).
\]
Omit the absent factors and terms when there are fewer children.
All references are in $(0,1)$ by induction. Integrating the
derivative of the log odds gives
$Q(\theta_b-d)\ge r_be^{-d}/(1-r_b+r_be^{-d})$.
Lowering the price by $d$ subtracts $dp_b$ from the scored
expression. Integrating the last bound and taking the supremum
over $d\ge0$ gives $\sup_d[-dp_b-\log(1-r_b+r_be^{-d})]
=K_-(p_b,r_b)$, proving \eqref{nu:reserve}. Every deduction is
spent at most once, at the parent.

\subsection{An integrated estimate at one split}

At the scoring vertex abbreviate $S=S_b$, $D=D(S\ell)$ and set
\[
 B=E(S\ell)-1-(E(S\ell)-\ell)u-\ell v,\qquad
 A=A_U(D),\qquad C_2=A_V(A).
\]
Here $B$ integrates the three bounds $E(S\ell),\ell,0$ for
zero, one and two active children, with the subtraction of one,
over the whole lifetime in units of $T_b$. The savings $A,C_2$
are obtained when the respective active children are fresh,
meaning none of their own children has yet split.

For a root at height one with child splits at $u\ge v>0$ and no
further splits, $D=\log(4/3)$ and the saving guaranteed by the
entropy bound relative to $B$ is $A(u-v)+C_2v$. As $u\downarrow0$,
the fresh reference values $q(U\ell)$ and $q(V\ell)$ tend to
$1/3$. Hence $A\to D/3$ and $C_2\to D/9$ by the definition of
$A_s$. The guaranteed saving is $D[(u-v)/3+v/9]+o(u)$.
Expanding $A(u-v)+C_2v$ gives the
coefficient $C_2-A\le0$ on $v$. Keeping the two interval lengths
$u-v$ and $v$ intact lets the certificate substitute lower bounds
for both savings directly.

For each child record its largest child split ratio $g$ and its
other ratio $z=gr$, with $0\le r\le1$ and $r=0$ if $g=0$.
A terminal child has $g=z=0$.
At a child with one active fresh descendant, its reference in
\eqref{nu:reserve} divided by its fresh reference is
$q(G(\ell-D(\theta)))$, where $G=g^{1-a}$ and $\theta$ is the
child's price. With two active fresh descendants this factor is
multiplied by $q(H(\ell-A_G(D(\theta))))$, $H=z^{1-a}$.
The factors are at least $1/3+0.09G$ and $1/3+0.124H$.
Indeed $q(t)\ge1/3+2t/9$ on $0\le t\le0.76\ell$ by
$q'=q(1-q)\ge2/9$. Also $\ell-D(\theta)\ge\log(3/2)>0.405$,
$D(\theta)<0.29$, and $A_G(d)\le q(G\ell)d\le0.46d$,
by concavity and $2^{0.76}<46/27$.

For $r_0\le q_s=q(s\ell)$ and $0\le d\le D$, writing
$z_0=q_s(1-e^{-sd})$ gives
\[
 A_{s,r_0}(d)\ge (r_0/q_s)\frac{2(1-z_0)}{2-z_0} A_s(d).
\]
This follows from $z\le-\log(1-z)\le z+z^2/[2(1-z)]$.
For $s\le0.76$ we have $z_0<0.092$, using
$0.76\cdot0.29<\log(5/4)$, hence the second factor exceeds
$0.951$. If $u\le h_0=0.225$ then $s\le u^{1-a}\le2/5$
by $(9/40)^{16}<(2/5)^{25}$ and $1-a>16/25$. In this case
$z_0\le0.46\cdot0.29\cdot2/5$ and the factor exceeds $0.97$.
The fractions for two children multiply because $A_{s,r}(d)$
is increasing, concave and zero at zero.

Use $f=0.97$ on parent $u$ intervals ending at or below $h_0$
and $f=0.951$ on the remaining interval. Let $p_1,p_2$ be the
cell constants below, satisfying $0\le p_2\le p_1\le0.36$,
$p_1\le f(1/3+0.09G)$ and
$p_2\le f(1/3+0.09G)(1/3+0.124H)$. Introduce an auxiliary
random cutoff for a child, in units of its own split height,
with distribution
\[
 \begin{array}{c|ccc}
 \text{value}&g&\max(z,h_0g)&h_0g\\
 \text{probability}&1-p_1&p_1-p_2&p_2 .
 \end{array}
\]
To justify the cutoff, first use the actual largest split ratio
$d_*$ among this child's grandchildren in place of $h_0g$, in
the same units. Above $g$ the child is fresh. Between
$\max(z,d_*)$ and $g$ it has one active fresh child, and down to
$d_*$ it has at most two active children, all fresh. Consequently
the probability that the time exceeds this cutoff is a lower
bound on its fraction of the fresh saving. Independent auxiliary
cutoffs on the two sides multiply these fractions.

Restoring $h_0g$ in the distribution overstates the fraction by
at most $0.36$ on an interval of absolute length bounded by
$\sum_c T_c(u_c-h_0)_+$, where $c$ ranges over internal children
of the child in question and $u_c$ is its largest child ratio.
Indeed $T_c\le gT_{\rm child}$ and $T_{\rm child}d_*=
\max_c T_cu_c$. Each $c$ has a single grandparent, the scoring
vertex $b$. Since $A,C_2\le0.46\cdot0.29$ and
$w_b/w_c\le0.76 S_c$, adding $0.037 S_c(u_c-h_0)_+$ to
the normalized cost at $c$ covers the error in the weighted sum.
Indeed $0.76\cdot0.36\cdot0.46\cdot0.29<0.037$. For fractions
$x,y,x',y'\in[0,1]$ one has
$(x'y'-xy)_+\le(x'-x)_++(y'-y)_+$, so the same charges cover
any overestimate of the product of the two saving fractions.

Let $X=uC_L$, $Y=vC_R$ for the two cutoffs in child units.
Integrating the one- and two-child saving fractions and using the
last correction now gives
\begin{align}
 \mathcal T_b={}& B-A\,\EE[u-\max(v,X)]
       -C_2\,\EE[(v-\max(X,Y))_+]+0.037S(u-h_0)_+,
     \notag\\[-2pt]
 &\hspace{30mm}0\le\sum_b T_b^a\mathcal T_b .                  \label{nu:toll}
\end{align}

\subsection{Rational certificate}

Set $J=(0.17,0.195,0.225)$ and $K=(0.67,0.8)$.
Divide the triangle $0\le z\le g\le h$ by the lines
$g\in\{0,0.135,0.17,0.195,0.225,h\}$ and
$z/g\in\{0,0.67,0.8,1\}$. In a cell take the lower endpoints
$g_0,r_0$ of the two scale ranges. For
$m(t)=10^{-3}\lfloor10^3t^{2/3}\rfloor$ and
$n(t)=10^{-4}\lfloor10^4 t\rfloor$, use
\[
 L_0=1/3+0.09m(g_0),\quad R_0=1/3+0.124m(g_0r_0),\qquad
 p_1=n(f L_0),\quad p_2=n(f L_0R_0).
\]
At vertices take the smaller of the adjacent lower endpoints
in each variable, or zero at the origin. These choices obey the
required bounds because $1-a<2/3$ and $g_0\le0.225$.

To choose a potential, consider the sum of
$T_w^a(\mathcal T_w+\varepsilon)$ over the internal vertices in
the subtree at $b$, divided by $T_b^a$, where
$\varepsilon=2\cdot10^{-6}$. This normalized sum equals
$\mathcal T_b+\varepsilon$ plus the two normalized child sums
with multipliers $u^a$ and $v^a$. We seek an upper bound
$F(S,u,v)$ that obeys the same recursion as an inequality.
The incoming price of vacancy is $S\ell$, and $u,v$ determine
both the weights of the two children and their prices $U\ell,V\ell$.
For fixed $S$ we allow the slope of $F$ to change on the lines
$u=J_i$ and $v=K_i u$. The resulting function is affine on each
scale cell, so its child contributions can be controlled at the
vertices. With $\xi(u,v)$ having first three entries $(u-J_i)_+$
and last two entries $(v-K_i u)_+$, take
\[
 P(S)=p_0+p'_1 S+p'_2 S^2,\qquad
 F(S,u,v)=P(S)-\sum_{i=1}^5 S(k_i+l_i S)\xi_i(u,v),
\]
with coefficients
\[
 \begin{aligned}
 10^7(p_0,p'_1,p'_2)&=(2998182,1908507,617448),\\
 10^7(k_1,\ldots,k_5)&=(-549378,518337,3850552,308836,2552524),\\
 10^7(l_1,\ldots,l_5)&=(2467210,-1612715,-5000000,-763799,-5000000).
 \end{aligned}
\]
For child scale pairs $(g,z)$ and $(g',z')$ the finite inequality is
\begin{equation}
 F(S,u,v)-\mathcal T_b-u^aF(U,g,z)-v^aF(V,g',z')
                      \ge 2\cdot10^{-6}.                      \label{nu:certificate}
\end{equation}
We specify its verification, implemented in the accompanying
\texttt{recursive\_certificate.py}.

For fixed cutoff values the portion of $\mathcal T_b$ before
the $0.037$ term is the maximum of
\[
 B-A(u-X),\quad B-A(u-v)-C_2(v-X),\quad
 B-A(u-v)-C_2(v-Y).
\]
This follows by comparing $X$ to $v$ and $Y$, using $0\le C_2\le A$.
The maximum is convex and nondecreasing in each cutoff. The cutoffs
are convex functions of the child scale pairs, with fixed probabilities
in each cell. The potential on a child is affine on its cell, so
it suffices to check \eqref{nu:certificate} at pairs of cell vertices.
Reducing both fractions at a vertex makes its cutoff stochastically
larger, which tests the more demanding inequality. For $u>0$, cut
$t=v/u$ at $0,0.67,0.8,1$, at possible values of $C_L$ and at
$C_L/C_R$ in $(0,1)$ for $C_R>0$. Both expected saving lengths
are rational linear combinations of $u,v$ on each resulting
interval. The lengths are nonnegative and the coefficients may
have either sign. Check both closed
pieces at shared boundaries and use continuity for $u=0$.

Here are the one-sided polynomial replacements used on each piece.
Write $u=x^{14}$, $t=y^{14}$, $S=0.76s$, with $s,x,y\in[0,1]$,
and $e=5-14a$. For $0\le z\le1$ put
\[
 W=z^{9+e},\quad w=z^9-ez^8(1-z^2)/2,\quad
 D_5=z^5+ez^4(1-z^2)/2+e^2z^3(1-z)^2/2,
 \quad D_{23}=z^{23}[1-e(1-z)-e(1-e)(1-z)^2/2].
\]
Then $w\le W\le z^9$, $W^2\ge z^{20}$,
$z^{5-e}\le D_5$ and $z^{23+e}\le D_{23}$.
The first and third bounds follow by expanding
$\exp(\pm e|\log z|)$ and using
$|\log z|\le (z^{-1}-z)/2$ and
$|\log z|\le(1-z)/\sqrt z$. For the positive exponential the
quadratic remainder is at most $e^2z^{-e}(\log z)^2/2$.
The last bound uses the Taylor expansion of $z^e$ at $1$
through degree two. Hence a branch potential $z^{5-e}P(W)$
is at most $p_0D_5+p'_1z^{14}+p'_2D_{23}$.
In the subtracted hinge charges $z^{5-e}W=z^{14}$ exactly.
For the remaining coefficient of $W$ choose $w$ for a lower
bound if it is nonnegative and $z^9$ otherwise.

Take rational brackets
$0.69314718055<\ell<0.69314718056$ and
$1.09861228866<\log3<1.09861228867$, using $-$ and $+$ for
the lower and upper endpoints. They and the three-row height
comparison follow from the first $120$ terms of the series for
$\log n=2\sum_{j\ge0} t^{2j+1}/(2j+1)$, $t=(n-1)/(n+1)$,
with tail at most $2t^{241}/[241(1-t^2)]$.
For $0\le b\le0.76\ell<0.527$ one has
$2/9\le q'(b)\le1/4$, $q''=q'(1-2q)$ and
$-1/8\le q'''=q'(1-6q')\le-2/27$. Taylor's formula gives
\[
 q(b)\ge1/3+2b/9+0.026b^2,\qquad E(S\ell)\le
 L=(\log3)_++\ell_+S/3+(\ell_+S)^2/9+
 (\ell_+S)^3/81-(\ell_-S)^4/324 .
\]
For the former inequality use $1/27-0.527/48>0.026$.
Thus $B\le L-1-(L-\ell_+)u-\ell_-v$ and
$0<d_0=2\ell_--L\le D<0.29$.
Integrating $q(Ub)$ over $[\ell-D,\ell]$ gives $A$.
The subinterval $[L-\ell_-,\ell_-]$ and the lower bound on $q$
give its minorant
\[
 A_-=d_0/3+\frac{w(x)}9[\ell_-^2-(L-\ell_-)^2]
        +\frac{0.026x^{20}}3[\ell_-^3-(L-\ell_-)^3].
\]
For the two-child term put
\[
 Q_z=1/3+2\ell_- w(z)/9+0.026\ell_-^2z^{20},\qquad
 G_z(d)=Q_z d-z^9 d^2(2/9+\ell_+z^9/12)/2 .
\]
We have $A_{W(z)}(d)\ge G_z(d)$ for $0\le d\le D$ by its
derivatives at zero and $q'(W\ell)\le2/9+W\ell/12$,
which follows from $q''\le1/12$. On this interval $G_z$ is
increasing and nonnegative since $Q_z>0.32$ and
$2/9+\ell_+/12<0.29$. Monotonicity then gives
$C_2\ge G_{xy}(G_x(d_0))$. The nonnegativity of the two
actual saving lengths justifies applying the minorants after
forming those lengths.

The verifier expands the rational polynomial minorants for the
slack minus $2\cdot10^{-6}$ in tensor-product Bernstein form.
In degree $d$ the contribution of the monomial $x^j$ to control
$i$ is $\binom{i}{j}/\binom{d}{j}$, understood as zero for
$i<j$. It encloses each component coefficient by rounding
down and up to multiples of $10^{-12}$. A signed multiplier uses
the corresponding lower or upper endpoint, and its product is
rounded down. Exact bisection on the mesh $2^{-25}$ encloses the
fourteenth roots of all domain endpoints. To restrict a Bernstein
table to $[L',H']$, its row $i$ averages the original controls
with weights given by the sum of independent
$\mathrm{Bin}(d-i,L')$ and $\mathrm{Bin}(i,H')$ variables,
as follows by expanding a Bernstein polynomial on that interval.
Subtracting the smallest control first makes all differences
nonnegative, so rounding each weight down to a multiple of
$2^{-45}$ and adding back the minimum gives a lower bound.
The program evaluates these averages using three base-$2^{15}$
digits, checking before each fixed-width matrix product that its
integer sums fit. Its initial component bounds have magnitude
below $2\cdot10^{12}$, the assembled sums use fewer than sixteen
such components of multiplier at most one in absolute value, and
midpoint subdivision only averages control values downwards.
Thus all subsequent signed additions stay in the 64-bit range.
If a coefficient is negative the box is subdivided at its midpoint,
cycling through $x,y,s$ and rounding every de Casteljau average
down. A nonnegative lower table certifies its whole box. The
complete counts of pieces, visited boxes and greatest depth are
\[
 \begin{array}{c|rrr}
 u\text{ interval}&\text{pieces}&\text{boxes}&\text{depth}\\ \hline
 {[0,0.17]}&3238&88490&31\\
 {[0.17,0.195]}&3238&6670&9\\
 {[0.195,0.225]}&3238&3414&5\\
 {[0.225,0.65]}&3238&4434&6 .
 \end{array}
\]
This gives \eqref{nu:certificate}. Multiplying it by $T_b^a$
and summing cancels all child potentials. By \eqref{nu:toll}
the resulting bound is
\[
 \sum_b T_b^a\le (2\cdot10^{-6})^{-1}p_0=C .
\]

\subsection{Improvement from a uniform bound on subtrees}

For each internal $v$ put $E_v=\sum_{w}T_w^a$, over all internal
descendants including $v$. Restricting to the rows of $v$ and
rescaling its positive lifetime to one gives another weak flag.
Indeed above $T_v$ all its columns are constant, so every local
test pulls back to an original test with full kernel there.
Hence $E_v\le C T_v^a$. For $\sigma=a_*/a$ we have
$2^{1-\sigma}(1-1/C)^\sigma=1$. Concavity now gives
\[
 \sum_{w\ {\rm internal\ child\ of}\ v} E_w^\sigma
 \le 2^{1-\sigma}(E_v-T_v^a)^\sigma\le E_v^\sigma.
\]
Every antichain of internal vertices therefore has total
$E_v^\sigma$ at most $C^\sigma$. The vertices with
$s\le T_v<s/h$ are an antichain because every child height is
at most $h$ times its parent height. Each such vertex contributes
at least $s^{a_*}$. Summing over $s=ch^{-j}$, $j\ge0$, bounds
the number of cells at height $c$ by
$1+C^\sigma c^{-a_*}/(1-h^{a_*})$. The variational formula for
$\alpha$ now gives \eqref{nu:bound}, by taking
$k=M(c)+1$, so $s_k\le c$, and letting $c\downarrow0$.
For the numerical comparison use $y=1/C$ and
$y+y^2/2\le-\log(1-y)\le y+y^2/2+y^3/[3(1-y)]$ together
with the rational bounds for $\ell$ above. This gives
$0.3556265775098<a_*<0.3556265775100$. Together with
Proposition~\ref{nl:result} this proves \eqref{e:decimal} and
completes Theorem~\ref{t:main}.

\appendix
\section{Realizing a strict flag on prime logarithms}
\label{s:realization}

A strict flag prescribes many ways to assign prime logarithms to
the rows. We need one assignment for which all the row sums are
close. Ford, Green and Koukoulopoulos first produce many attainable
discrepancy vectors in the integer model and then use fresh elements
to reach the diagonal \cite[Propositions~5.9, 6.4 and Lemma~6.5]{fgk}.
They trace the latter step to Maier and Tenenbaum \cite{mt84}.
Mao and Song \cite[Appendix~A]{ms} give corrections to the moment
and lattice estimates in the integer argument.

For prime logarithms the attainable vectors lie in a real box.
We use a second moment to show that unit neighbourhoods of these
vectors cover a positive proportion of the box. A fresh pair of
primes assigned complementary columns then moves one discrepancy
coordinate by $\log(p/q)$. One such pair per coordinate yields a
positive probability of reaching bounded distance from the
diagonal. Conditional on the process supplying the attainable
vectors for all trials, the trials use disjoint, independent
reserves.

As in Section~\ref{s:arith}, $\mathcal P$ has intensity
$\nu=\sum_p p^{-1}\delta_{\log p}$, denoted there by $\mu$.
Lemma~\ref{l:prime-log-measure} gives
$\sum_{\log p\le y}(\log p)/p=y+O(1)$.
Thus, for an absolute $C>0$, every half-open interval with endpoints
$0<a<d$ satisfies
\begin{equation}
 \frac{d-a-C}{d}\le\nu((a,d])\le\frac{d-a+C}{a},\qquad
 \nu((a,d])=\log(d/a)+O(1/a).                              \label{e:real-intensity}
\end{equation}
The first bounds follow by bounding $1/y$ on the interval, and the
last equality follows by partial summation. Endpoint conventions are
immaterial after enlarging $C$.

\begin{lemma}\label{l:realization}
Suppose a separating Boolean flag on $k\ge2$ rows has rank-one births
at $1\ge c_1>\cdots>c_r>c_{r+1}=b>0$. Suppose that its stage laws
satisfy the strict positive-endpoint tests. For each $0<b'<b$, with
probability tending to one as $s\to\infty$, there are $k$ distinct
subsets of $\mathcal P\cap(\exp(b's),\exp s]$ whose sums have
diameter bounded by a constant depending only on the flag. The primes
underlying all the selected points may be taken distinct.
\end{lemma}

\begin{proof}
Work modulo the line of constant row vectors. Identify the resulting
terminal space with $\mathbb R^r$ using the images of the successive
Boolean birth columns $u_1,\dots,u_r$ as basis. Write $V_j$ for the
span of the first $j$ basis vectors in these coordinates, $V_0=0$,
and $\mu_j$ for the image of the stage law. All its letters have
bounded coordinates, with a bound depending only on this fixed flag.
Every assignment of letters below uses Boolean representatives from
the original stage alphabets. Its sum in $\mathbb R^r$ specifies the
row sums modulo their common translation.

There is $g>0$ such that each proper boundary test has excess at
least $g$. Indeed the dimensions of its kernels and the partitions
they induce on the finite stage alphabets have only finitely many
profiles. We need a bound also for kernels which grow inside a stage.
Use intervals of lengths $s(c_j-c_{j+1})$ in descending order, and
write $W$ for any increasing family of rational kernels with
$W\subseteq V_j$ on the $j$th interval. Let $D$ be the total length
on which $W\ne V_j$, and $e$ the final codimension. Sampling a uniform
time in each interval gives a boundary test. Its probability of being
proper is at least $D/s$, since this is bounded by the largest
fraction of any interval spent with a proper kernel. Its sampled
terminal codimension is at least $e$. Averaging the boundary tests
therefore gives
\[
 \sum_j\int_j [H(\mu_j\bmod W)-\dim(V_j/W)]\,dt-bs e\ \ge\ gD.
\]
Here $\int_j$ denotes integration over the $j$th interval in the
physical logarithmic coordinate $t=\log y$. Its origin does not
affect this inequality. The entropy vanishes whenever $W=V_j$ and
is bounded on every stage by $H_0=\max_j\log|\operatorname{supp}\mu_j|$.
Fix $0<\lambda<1$ sufficiently close to one that
$(1-\lambda)H_0<g/2$. We obtain
\begin{equation}
 \sum_j\int_j [\lambda H(\mu_j\bmod W)-\dim(V_j/W)]\,dt-bs e
       \ \ge\ (g/2)D .                                      \label{e:real-gap}
\end{equation}

We now construct the attainable set. Independently color the process
into three Poisson processes of
intensities $\lambda\nu$ and $(1-\lambda)\nu/2$ twice. The
first supplies a set of possible sums. For a translation $h\ge0$
write $a_j=c_js-h$, $1\le j\le r+1$, and $U_j=\exp a_j$.
In the bulk estimate allow $0\le h\le bs/2$, so $a_{r+1}\ge bs/2$.
In each logarithmic interval $(a_{j+1},a_j]$, use consecutive bins of
a fixed length $L\ge1$ starting from the top. Discard any final
partial bin. All constants below may depend on the flag. The value
of $L$ will be fixed independently of $s,h$.

In a complete bin of stage $j$ with $n$ points of the first color,
choose nonnegative integer counts for the letters, summing to $n$
and differing from $n\mu_j$ by at most one per letter. Use a uniform
random permutation of that multiset on the points, independently in
the different bins. This defines a probability law on possible
discrepancy vectors $Z=\sum_y y\,u(y)$ for each fixed configuration
of points. We claim that for every $\xi>0$ there
are $A,C_1<\infty$ such that, uniformly in the indicated translations,
with probability at least $1-\xi$ for all sufficiently large $s$,
\begin{align}
 &\operatorname{supp} Z\ \subseteq\ \prod_{j=1}^r[-A U_j,A U_j],
       \nonumber\\
 &\left|\operatorname{supp} Z+[-1,1]^r\right|
          \ \ge\ C_1^{-1}\prod_{j=1}^r U_j .                   \label{e:real-volume}
\end{align}
The bars in the second line denote Lebesgue measure. Coordinate $j$
gets contributions only from points below $U_j$. The expected sum of
all first-color points up to $U_j$ is $\lambda(U_j+O(1))$. The coordinate
bound on the letters and Markov's inequality therefore prove the
first line outside an event of arbitrarily small probability.

To prove the second line, fix a small $\delta>0$. At each stage start
require that the number $N(v)$ of first-color points in the first $v$
logarithmic units satisfy
\begin{equation}
              |N(v)-\lambda v|\le\delta v+C_0                 \label{e:real-count}
\end{equation}
at all complete bin boundaries. The probability of failure is
arbitrarily small when the constant $C_0$ is large, uniformly in
$s,h$. For example, choose $\theta>0$ small enough that both
$\lambda(e^{\pm\theta}-1\mp\theta)<\theta\delta/2$. At every tested
prefix the mean is $\lambda v+O(1)$ by \eqref{e:real-intensity},
uniformly since all used locations tend to infinity. The Poisson
exponential martingales, with these deterministic cumulative means,
and their maximal inequality bound the probability of failure on
each side by $O(e^{-\theta C_0})$. Sum over the fixed number of stage starts.

Condition on bin counts satisfying \eqref{e:real-count}. Label the
points within a bin in uniformly random order. Their locations are
then independent draws from $\nu$ normalized on that bin. Assign two
independent letter permutations to these labels, giving vectors $Z,Z'$.
The permutations are independent of the locations conditional on the
counts. For a pair of permutations, take the cumulative span of its
letter differences as bins are read from the top down. Record each
bin in which this span increases and the resulting subspace. There
are at most $r$ such bins. Define $W$ using the updated span from the
upper boundary of each recorded bin, with $W=0$ initially. On discarded
partial bins the span stays fixed. This gives nested kernels with
$W\subseteq V_j$. Let $r'=\dim W_{\rm end}$, $e=r-r'$, and list
the upper logarithmic bin endpoints $a'_1,\ldots,a'_{r'}$ at which
dimensions were added, repeating each endpoint by its increment.
The rank identity is
\begin{equation}
 \sum_{j=1}^r a_j-\sum_{i=1}^{r'}a'_i
       =\sum_j\int_j\dim(V_j/W)\,dt+a_{r+1} e.                \label{e:real-rank}
\end{equation}
It follows by integrating the dimension contributions from each birth
down to the end of the range. Since $a_{r+1}\le bs$, the last term
is covered by \eqref{e:real-gap}.

For a fixed pair of permutations, choose $r'$ independent difference
letters from their recorded bins, adding a basis of each increment
in turn. Condition on all other locations. The chosen locations are
independent. A nonzero $r'$-row minor of the chosen columns confines
each of these locations to an interval of length $O(1)$ when
$\|Z-Z'\|_\infty\le2$. The constant is uniform since the columns come
from a fixed finite alphabet. For a bin with upper endpoint $a'$, its
total $\nu$-mass is $L+o(1)$, while any interval of bounded length
inside it has $\nu$-mass $O_L(e^{-a'})$ by \eqref{e:real-intensity}.
Hence the conditional point lands in its required interval with
probability $O_L(e^{-a'})$. For
$r'=0$ use probability at most one. Thus the conditional small-ball
probability is at most
\begin{equation}
            O_L(1)\exp\left(-\sum_i a'_i\right).
                                                               \label{e:real-pivots}
\end{equation}

Now fix the recorded bins and subspaces, which we call a profile in
this proof. In an unrecorded bin of stage $j$ with current subspace $W$,
the two permutations must agree point by point modulo $W$. The
probability of agreement is the reciprocal of the multinomial
coefficient for the bin's quotient counts. Hence its logarithm is
at most $-nH(\mu_j\bmod W)+O(\log(n+2))$. Indeed, the factorial
estimate $\log(n!)=n\log n-n+O(\log(n+2))$ and the bounded rounding
errors give exactly this error for each of the finitely many letters
and quotient classes. Use the estimate only for $W\ne V_j$, and
impose no constraints on recorded bins. Independence over bins bounds
the logarithm of the profile probability by
\begin{equation}
 -\sum_j\int_j\lambda H(\mu_j\bmod W)\,dt
   +O\left((\delta+\log(L+2)/L)D\right)+O_{C_0,L}(1).
                                                               \label{e:real-profile}
\end{equation}
Here the coefficient in the first error depends only on the fixed
flag. We detail the estimate because a bound growing with the total
length of the clock would not suffice. In each stage the times with
$W\ne V_j$ form an initial segment of length $D_j$, and $\sum_jD_j=D$.
The unrecorded complete bins with a fixed proper $W$ form runs, at
most $2r+1$ over all stages. Both endpoints of any such run are
within $D_j+L$ of its stage start. Subtracting the two prefix bounds
\eqref{e:real-count} replaces its count by $\lambda$ times its length
at an error at most $2\delta(D_j+L)+2C_0$. The omitted transition
and partial bins have total length at most $2rL$. In a stage with $m$
complete bins in its deficient prefix their total count is bounded
by $(\lambda+\delta)mL+C_0$. Concavity of the logarithm bounds the
sum of the multinomial errors there by a constant times
$m\log(2+(\lambda+\delta)L)+C_0$. This proves \eqref{e:real-profile}.

Take $L$ large, then $\delta$ small, so that the coefficient times
$D$ in \eqref{e:real-profile} is at most $gD/8$. Combine the drift
\eqref{e:real-gap} with the rank identity \eqref{e:real-rank} and the
two probability bounds \eqref{e:real-pivots} and \eqref{e:real-profile}. Relative to
$\exp(-\sum_j a_j)$, the contribution of any fixed profile
to the small-ball probability is bounded by
$O_{C_0,L}(1)e^{-3gD/8}$.
For $m\le D<m+1$ the number of profiles is $O((m+2)^r)$ with a
constant independent of $s,h$. Every recorded bin begins in, or at
the end of, its stage's deficient prefix and is thus among its first
$1+(m+1)/L$ bins. The stages, increments and possible subspaces are
drawn from fixed finite lists, since every span is generated by
letter differences. Summing the convergent series gives
\[
 \Pr(\|Z-Z'\|_\infty\le2\mid\hbox{counts})
        \ll_{C_0,L} \left(\prod_j U_j\right)^{-1}.
\]
This probability averages over both positions and permutations.
For fixed positions smooth the letter-assignment law by the uniform
law on $[-1,1]^r$ and call the density $f$. Its squared integral
is bounded by $2^{-r}$ times the preceding small-ball probability
conditional on the positions, as follows by integrating the overlap
of two translated cubes. Markov's inequality now gives
$\int f^2\le C_1/\prod_j U_j$ outside an event of probability as
small as desired over the first-color configuration. By
Cauchy--Schwarz its support has volume at least $1/\int f^2$.
This proves \eqref{e:real-volume}.

It remains to turn positive volume into a close family of divisors.
Fix $\xi>0$ and the corresponding
constants in \eqref{e:real-volume}.
Put $K=10(A+2)$ and fix $B>\log K$. Use $m_s=\lfloor\sqrt{s}\rfloor$
trials with translations $h=lB$, $1\le l\le m_s$. For each direction
$j$ and trial use the reserve interval $[U_j,KU_j]$. For large $s$
these intervals are disjoint over all pairs $(j,l)$. This follows
from $B>\log K$ for a common $j$ and from the fixed positive gaps
between successive $c_j$ for distinct directions. Every used point
lies in $(\exp(b's),\exp s]$, since $c_1\le1$, the first translation
exceeds $\log K$, and $(b-b')s>Bm_s$ eventually.

Condition on the first-color configuration. For one trial, require
exactly one point of each of the two other colors in each of its $r$
reserve intervals. This event has probability bounded below by a
positive constant uniformly in $s,l$, since each interval has mass
$\log K+o(1)$. The two points in direction $j$ have independent laws
given by $\nu$ normalized on $[U_j,KU_j]$. Fix $H>C+2$ with $C$ as
in \eqref{e:real-intensity}. For every half-open interval $I$ of
length $H$ contained in $[-(A+2)U_j,(A+2)U_j]$, the difference
of these points lies in $I$ with probability $\gg_A 1/U_j$.
Indeed the first point lies in $[(A+4)U_j,(K-A-4)U_j]$ with
probability bounded below by \eqref{e:real-intensity}. Given such a
point $y$, the interval $y-I$ is inside $[U_j,KU_j]$ and has mass
at least $(H-C)/(KU_j)$. Assign the Boolean column
$u_j$ to the first point and its complement to the other. Their net
contribution in the chosen basis is this difference in coordinate $j$.
The contributions in different coordinates are independent. Use the
disjoint half-open cubes of side $H$ on the standard grid in
$\mathbb R^r$. At least $(C_1H^r)^{-1}\prod_j U_j$ such cubes meet
the negative of the smoothed support in \eqref{e:real-volume}. Every
one of these cubes is in $\prod_j[-(A+2)U_j,(A+2)U_j]$ for large $s$.
The preceding interval estimate gives probability $\gg_A1/\prod_j U_j$
to each cube. On their union some bulk assignment brings all
discrepancy coordinates within $H+1$ of zero. The conditional
probability of success in a trial satisfying \eqref{e:real-volume}
is therefore at least $p_\xi>0$ independent of $s,l$.

The expected fraction of trials failing \eqref{e:real-volume} is
at most $\xi$. With probability at least $1-2\xi$ at least half the
trials satisfy it, by Markov's inequality. Their reserve events are
conditionally independent because their intervals are disjoint and
their colors independent of the first process. The probability that
all trials fail is at most $2\xi+(1-p_\xi)^{m_s/2}$. Let $s\to\infty$
and then $\xi\downarrow0$. For a successful trial, every row pair is
distinguished by at least one assigned reserve column because the
columns $u_j$ separate the rows. Each row-sum difference has absolute
value at most $r(H+1)$, since all entries of each $u_j$ are zero or
one. The probability of two occurrences of any prime in the entire
range is at most $\tfrac12\sum_{p\ge\exp(\exp(b's))}p^{-2}=o(1)$.
The diameter bound $r(H+1)$ is independent of the chosen failure
tolerance $\xi$. This proves the lemma.
\end{proof}

AI Disclosure: The author selected the research topic; artificial intelligence completed all mathematical derivations, proofs, and writing.

\begin{thebibliography}{99}

\bibitem{bt}
R. de la Bret\`eche and G. Tenenbaum,
Two upper bounds for the Erd\H{o}s--Hooley Delta-function,
\emph{Sci. China Math.} \textbf{66} (2023), 2683--2692.
\href{https://doi.org/10.1007/s11425-022-2193-8}{doi:10.1007/s11425-022-2193-8}.

\bibitem{ct}
T. M. Cover and J. A. Thomas, \emph{Elements of Information Theory},
second edition, Wiley-Interscience, Hoboken, NJ, 2006.

\bibitem{erdos}
P. Erd\H{o}s, On the density of some sequences of integers,
\emph{Bull. Amer. Math. Soc.} \textbf{54} (1948), 685--692.
\href{https://doi.org/10.1090/S0002-9904-1948-09057-7}{doi:10.1090/S0002-9904-1948-09057-7}.

\bibitem{en75}
P. Erd\H{o}s and J.-L. Nicolas, R\'epartition des nombres superabondants,
\emph{Bull. Soc. Math. France} \textbf{103} (1975), 65--90.
\href{https://doi.org/10.24033/bsmf.1793}{doi:10.24033/bsmf.1793}.

\bibitem{en76}
P. Erd\H{o}s and J.-L. Nicolas, M\'ethodes probabilistes et combinatoires
en th\'eorie des nombres, \emph{Bull. Sci. Math.} (2) \textbf{100}
(1976), \href{https://www.renyi.hu/~p_erdos/1976-16.pdf}{301--320}.

\bibitem{fgk}
K. Ford, B. Green and D. Koukoulopoulos,
Equal sums in random sets and the concentration of divisors,
\emph{Invent. Math.} \textbf{232} (2023), 1027--1160.
\href{https://doi.org/10.1007/s00222-022-01177-y}{doi:10.1007/s00222-022-01177-y}.

\bibitem{ht}
R. R. Hall and G. Tenenbaum,
On the average and normal orders of Hooley's $\Delta$-function,
\emph{J. London Math. Soc.} (2) \textbf{25} (1982), 392--406.
\href{https://doi.org/10.1112/jlms/s2-25.3.392}{doi:10.1112/jlms/s2-25.3.392}.

\bibitem{hooley}
C. Hooley, On a new technique and its applications to the theory of numbers,
\emph{Proc. London Math. Soc.} (3) \textbf{38} (1979), 115--151.
\href{https://doi.org/10.1112/plms/s3-38.1.115}{doi:10.1112/plms/s3-38.1.115}.

\bibitem{hz}
J. Hou and H. Zhao, Equal subset sums and close divisors,
preprint (2026), \href{https://arxiv.org/abs/2609.26014}{arXiv:2609.26014}.

\bibitem{mt84}
H. Maier and G. Tenenbaum, On the set of divisors of an integer,
\emph{Invent. Math.} \textbf{76} (1984), 121--128.
\href{https://doi.org/10.1007/BF01388495}{doi:10.1007/BF01388495}.

\bibitem{mt85}
H. Maier and G. Tenenbaum, On the normal concentration of divisors,
\emph{J. London Math. Soc.} (2) \textbf{31} (1985), 393--400.
\href{https://doi.org/10.1112/jlms/s2-31.3.393}{doi:10.1112/jlms/s2-31.3.393}.

\bibitem{mt09}
H. Maier and G. Tenenbaum, On the normal concentration of divisors, 2,
\emph{Math. Proc. Cambridge Philos. Soc.} \textbf{147} (2009), 513--540.
\href{https://doi.org/10.1017/S0305004109002631}{doi:10.1017/S0305004109002631}.

\bibitem{ms}
Y. Mao and Y. Song, Close divisors of typical integers:
the Ford--Green--Koukoulopoulos conjecture,
preprint (2026), \href{https://arxiv.org/abs/2609.22296}{arXiv:2609.22296}.

\end{thebibliography}
\end{document}
